% Generated mechanically from the frozen relative-cycles.tex target.
% Do not edit this generated file; edit the bound locale source instead.

% OK, start here.
%

\setcounter{chapter}{61}
\chapter{相对循环}



\phantomsection
\label{relative-cycles-section-phantom}




\section{引言}
\label{relative-cycles-section-introduction}

\noindent
基础参考文献是 \cite{SV}。

\medskip\noindent
本章只定义 \cite{SV} 中所谓的普遍整相对循环。与原文相比，这一选择使理论的
展开略为简洁，但当然也会损失一些内容。

\medskip\noindent
固定 Noether 概形之间的有限型态射 $X \to S$。$X/S$ 各纤维上的
$r$-循环族 $\alpha$，就是一个集合 $\alpha = (\alpha_s)_{s \in S}$，
其中 $\alpha_s \in Z_r(X_s)$。沿任意态射 $g : S' \to S$ 对 $\alpha$
作基变换而得到 $g^*\alpha$ 的方式是显然的。若 $\alpha$ 与特化相容，
则称 $\alpha$ 为 {\it $X/S$ 上的相对 $r$-循环}。换言之，对任意态射
$g : S' \to S$，其中 $S'$ 是一个离散赋值环的谱，我们要求
$g^*\alpha$ 的一般纤维特化为 $g^*\alpha$ 的闭纤维。见第
\ref{relative-cycles-section-families-specialization} 节。






\section{约定与记号}
\label{relative-cycles-section-conventions}

\noindent
关于固定 Noether 基上局部有限型概形之循环的约定与记号，请参阅
“Chow 同调与陈类”一章；见 Chow 同调，第 \ref{chow-section-setup} 节及后文。

\medskip\noindent
特别地，若 $X$ 在域 $k$ 上局部有限型，则 $Z_r(X)$ 表示 $r$ 维循环群；
见 Chow 同调，例 \ref{chow-example-field} 与第
\ref{chow-section-cycles} 节。给定满足 $\dim(Z) = r$ 的整闭子概形
$Z \subset X$，有 $[Z] \in Z_r(X)$；若 $X$ 拟紧，则 $Z_r(X)$ 是以这些
类为基的自由阿贝尔群。




\section{相对于域的循环}
\label{relative-cycles-section-relative-fields}

\noindent
设 $k$ 为域，$X$ 为 $k$ 上的局部代数概形，并设 $r \geq 0$ 为整数。
在此情形下，我们有 $X$ 上的 $r$-循环群 $Z_r(X)$；见第
\ref{relative-cycles-section-conventions} 节。

\medskip\noindent
{\bf 基变换。} 对任意域扩张 $k'/k$，都有基变换映射
$Z_r(X) \to Z_r(X_{k'})$；见 Chow 同调，第
\ref{chow-section-change-base} 节。具体而言，给定维数为 $r$ 的整闭子概形
$Z \subset X$，我们把 $[Z] \in Z_r(X)$ 映到与闭子概形
$Z_{k'} \subset X_{k'}$ 相伴的 $r$-循环
$[Z_{k'}]_r \in Z_r(X_{k'})$（一般而言，$Z_{k'}$ 当然既不必不可约，
也不必既约）。基变换映射 $Z_r(X) \to Z_r(X_{k'})$ 总是单射。

\begin{lemma}
\label{relative-cycles-lemma-multiplicities-field-extension}
设 $K/k$ 为域扩张，$Z$ 为 $k$ 上的整局部代数概形。不可约分支
$Z' \subset Z_K$ 的重数 $m_{Z', Z_K}$ 等于 $1$ 或 $k$ 的特征的某个幂。
\end{lemma}

\begin{proof}
若 $k$ 的特征为零，则 $k$ 完美，并且由“代数簇”引理
\ref{varieties-lemma-geometrically-reduced}，$X_K$ 既约，故重数总是 $1$。
现假设 $k$ 的特征为 $p > 0$。设 $L$ 为 $Z$ 的函数域。由于 $Z$ 在 $k$
上局部代数，域扩张 $L/k$ 有限生成。环 $K \otimes_k L$ 是 Noether 环
（代数，引理 \ref{algebra-lemma-Noetherian-field-extension}）。转化为代数
语言，我们必须证明：对每个极小素理想 $\mathfrak q$，Artin 局部环
$(K \otimes_k L)_\mathfrak q$ 的长度都是 $p$ 的幂。

\medskip\noindent
设 $L'/L$ 为有限纯不可分扩张，例如次数为 $p^n$。则
$K \otimes_k L \subset K \otimes_k L'$ 是秩为 $p^n$ 的有限自由环映射；
它在谱上诱导同胚，并诱导纯不可分的剩余域扩张。因此，对上述每个极小素
理想 $\mathfrak q$，都有唯一的极小素理想
$\mathfrak q' \subset K \otimes_k L'$ 位于其上，并且
$$
p^n \text{length}((K \otimes_k L)_\mathfrak q) =
[\kappa(\mathfrak q') : \kappa(\mathfrak q)]
\text{length}((K \otimes_k L')_{\mathfrak q'})
$$
这是把代数，引理 \ref{algebra-lemma-pushdown-module} 应用于
$M = (K \otimes_k L')_{\mathfrak q'} \cong
(K \otimes_k L)_{\mathfrak q}^{\oplus p^n}$ 所得。由于
$[\kappa(\mathfrak q') : \kappa(\mathfrak q)]$ 是 $p$ 的幂，可知只需对
$L'$ 与 $\mathfrak q'$ 证明该断言。

\medskip\noindent
由前一段和代数，引理 \ref{algebra-lemma-make-separable}，可以假设存在
子域塔 $L/k'/k$，使得 $L/k'$ 可分且 $k'/k$ 有限纯不可分。于是
$K \otimes_k k'$ 是 Artin 局部环。前一段的论证（取 $L = k$、
$L' = k'$）表明 $\text{length}(K \otimes_k k')$ 是 $p$ 的幂。又由于
$L/k'$ 是一个光滑 $k'$-代数的局部化（代数，引理
\ref{algebra-lemma-localization-smooth-separable}），故
$S = (K \otimes_k L)_\mathfrak q$ 是某个光滑
$R = K \otimes_k k'$-代数在极小素理想处的局部化。因此
$R \to S$ 是 Artin 局部环之间的平坦局部同态，且
$\mathfrak m_R S = \mathfrak m_S$。由代数，引理
\ref{algebra-lemma-pullback-module} 得
$\text{length}(K \otimes_k k') = \text{length}(R) =
\text{length}(S) = \text{length}((K \otimes_k L)_\mathfrak q)$
从而证明完成。
\end{proof}

\begin{lemma}
\label{relative-cycles-lemma-how-different}
设 $k$ 为特征 $p > 0$ 的域，其完美闭包为 $k^{perf}$。设 $X$ 为 $k$ 上的
代数概形，$r \geq 0$ 为整数。单射
$Z_r(X) \to Z_r(X_{k^{perf}})$ 的余核是 $p$-幂挠模（代数续篇，定义
\ref{more-algebra-definition-f-power-torsion}）。
\end{lemma}

\begin{proof}
由于 $X$ 拟紧，阿贝尔群 $Z_r(X)$ 是自由的，其基由维数为 $r$ 的整闭子概形
给出；$Z_r(X_{k^{perf}})$ 亦然。由于
$X_{k^{perf}} \to X$ 是同胚，映射 $Z_r(X) \to Z_r(X_{k^{perf}})$
为单射且余核为挠群。由引理 \ref{relative-cycles-lemma-multiplicities-field-extension}，
余核中的每个元素都是 $p$-幂挠的。
\end{proof}





\section{循环的特化}
\label{relative-cycles-section-specialization}

\noindent
设 $R$ 为离散赋值环，其分式域为 $K$，剩余域为 $\kappa$。设 $X$ 为
$R$ 上局部有限型的概形，并设 $r \geq 0$。存在特化映射
$$
sp_{X/R} : Z_r(X_K) \longrightarrow Z_r(X_\kappa)
$$
其定义如下。对维数为 $r$ 的整闭子概形 $Z \subset X_K$，以
$\overline{Z}$ 表示 $Z \to X$ 的概形论像。令 $sp_{X/R}$ 为满足下式的
唯一 $\mathbf{Z}$-线性映射：
$$
sp_{X/R}([Z]) = [\overline{Z}_\kappa]_r
$$
简要说明该定义为何良定。首先，态射 $X_K \to X$ 拟紧，故态射
$Z \to X$ 拟紧。因此，由“态射”引理
\ref{morphisms-lemma-flat-base-change-scheme-theoretic-image}，取
$Z \to X$ 的概形论像与平坦基变换可交换。特别地，基变换回 $X_K$ 后，
有 $Z = \overline{Z}_K$。由于 $Z$ 是整的，$\overline{Z}$ 当然也是整的；
事实上，它等于以 $Z$ 的一般点（之像）为一般点的唯一整闭子概形。由
“代数簇”引理
\ref{varieties-lemma-dominate-valuation-ring-dimension-fibres}，
$Z_\kappa$ 是维数为 $r$ 的等维概形。

\begin{lemma}
\label{relative-cycles-lemma-specialization-module}
设 $R$ 为离散赋值环，其分式域为 $K$，剩余域为 $\kappa$。设 $X$ 为
$R$ 上局部有限型的概形，$r \geq 0$。设 $\mathcal{F}$ 为在 $R$ 上平坦的
凝聚 $\mathcal{O}_X$-模。假设
$\dim(\text{Supp}(\mathcal{F}_K)) \leq r$。则
$\dim(\text{Supp}(\mathcal{F}_\kappa)) \leq r$，并且
$$
sp_{X/R}([\mathcal{F}_K]_r) = [\mathcal{F}_\kappa]_r
$$
\end{lemma}

\begin{proof}
关于维数的断言由“态射续篇”引理
\ref{more-morphisms-lemma-relative-dimension-support-flat} 得出。
设 $x$ 为维数 $r$ 的整闭子概形 $Z \subset X_\kappa$ 的一般点。为完成
证明，我们将说明等式左边 (L) 与右边 (R) 中 $[Z]$ 的系数相同。

\medskip\noindent
令 $A = \mathcal{O}_{X, x}$、$M = \mathcal{F}_x$。注意，$M$ 是在
$R$ 上平坦的有限 $A$-模。取素元 $\pi \in R$，使得
$A/\pi A = \mathcal{O}_{X_\kappa, x}$。由 Chow 同调，引理
\ref{chow-lemma-additivity-divisors-restricted}，有
$$
\sum\nolimits_i \text{length}_A(A/(\pi, \mathfrak q_i))
\text{length}_{A_{\mathfrak q_i}}(M_{\mathfrak q_i}) =
\text{length}_A(M/\pi M)
$$
其中求和遍历 $M$ 的支撑中的极小素理想 $\mathfrak q_i$。由于 $\pi$ 是
$M$ 上的非零因子，有 $\pi \not \in \mathfrak q_i$；因此，这些素理想
对应于 $\mathcal{F}_K$ 的支撑中那些特化到所选
$x \in X_\kappa$ 的一般点 $y_i \in X_K$。所以左边就是 (L) 中 $[Z]$
的系数。当然，$\text{length}_A(M/\pi M)$ 是 (R) 中 $[Z]$ 的系数，
证明完成。
\end{proof}

\begin{lemma}
\label{relative-cycles-lemma-specialization-closed}
设 $R$ 为离散赋值环，其分式域为 $K$，剩余域为 $\kappa$。设 $X$ 为
$R$ 上局部有限型的概形，$r \geq 0$。设 $W \subset X$ 为在 $R$ 上平坦的
闭子概形。假设 $\dim(W_K) \leq r$。则 $\dim(W_\kappa) \leq r$，并且
$$
sp_{X/R}([W_K]_r) = [W_\kappa]_r
$$
\end{lemma}

\begin{proof}
取 $\mathcal{F} = \mathcal{O}_W$，这就是引理
\ref{relative-cycles-lemma-specialization-module} 的特殊情形。见 Chow 同调，引理
\ref{chow-lemma-cycle-closed-coherent}。
\end{proof}

\begin{lemma}
\label{relative-cycles-lemma-specialization-extension}
设 $R'/R$ 为离散赋值环的扩张，它诱导分式域扩张 $K'/K$ 与剩余域扩张
$\kappa'/\kappa$（代数续篇，定义
\ref{more-algebra-definition-extension-discrete-valuation-rings}）。
设 $X$ 在 $R$ 上局部有限型，并记 $X' = X_{R'}$。则图
$$
\xymatrix{
Z_r(X'_{K'}) \ar[rr]_{sp_{X'/R'}} & & Z_r(X'_{\kappa'}) \\
Z_r(X_K) \ar[rr]^{sp_{X/R}} \ar[u] & & Z_r(X_\kappa) \ar[u]
}
$$
可交换，其中 $r \geq 0$，竖直箭头是基变换映射。
\end{lemma}

\begin{proof}
注意，$X'_{K'} = X_{K'} = X_K \times_{\Spec(K)} \Spec(K')$；闭纤维亦然，
故竖直箭头确有意义（见第 \ref{relative-cycles-section-relative-fields} 节）。现设
$Z \subset X_K$ 为整闭子概形，其概形论像为 $\overline{Z} \subset X$。
则 $Z_{K'} \subset X_{K'}$ 是闭子概形，其概形论像为
$\overline{Z}_{R'} \subset X_{R'}$。按定义，$[Z]$ 的基变换是
$[Z_{K'}]_r = [\overline{Z}_{K'}]_r$。我们有
$$
sp_{X/R}([Z]) = [\overline{Z}_\kappa]_r
\quad\text{且}\quad
sp_{X'/R'}([\overline{Z}_{K'}]_r) = [(\overline{Z}_{R'})_{\kappa'}]_r
$$
这是由引理 \ref{relative-cycles-lemma-specialization-module} 得出的。又因为
$(\overline{Z}_{R'})_{\kappa'} = (\overline{Z}_\kappa)_{\kappa'}$，
故得结论。
\end{proof}

\begin{lemma}
\label{relative-cycles-lemma-specialization-flat-pullback}
设 $R$ 为离散赋值环，其分式域为 $K$，剩余域为 $\kappa$。设 $X$ 为
$R$ 上局部有限型的概形。设 $f : X' \to X$ 为局部有限型、平坦且相对
维数为 $e$ 的态射。则图
$$
\xymatrix{
Z_{r + e}(X'_K) \ar[rr]_{sp_{X'/R}} & & Z_{r + e}(X'_\kappa) \\
Z_r(X_K) \ar[rr]^{sp_{X/R}} \ar[u] & & Z_r(X_\kappa) \ar[u]
}
$$
可交换，其中 $r \geq 0$，竖直箭头由平坦拉回给出。
\end{lemma}

\begin{proof}
设 $Z \subset X$ 为支配 $R$ 的整闭子概形。由 $sp_{X/R}$ 的构造，有
$sp_{X/R}([Z_K]) = [Z_\kappa]_r$，而这刻画了特化映射。置
$Z' = f^{-1}(Z) = X' \times_X Z$。由于 $R$ 是赋值环，$Z$ 在 $R$ 上
平坦。因此 $Z'$ 在 $R$ 上平坦，并且由引理
\ref{relative-cycles-lemma-specialization-closed}，有
$sp_{X'/R}([Z'_K]_{r + e}) = [Z'_\kappa]_{r + e}$
又由 Chow 同调，引理 \ref{chow-lemma-pullback-coherent}，有
$f_K^*[Z_K] = [Z'_K]_{r + e}$ 及
$f_\kappa^*[Z_\kappa]_r = [Z'_\kappa]_{r + e}$，故得结论。
\end{proof}

\begin{lemma}
\label{relative-cycles-lemma-specialization-proper-pushforward}
设 $R$ 为离散赋值环，其分式域为 $K$，剩余域为 $\kappa$。设
$f : X \to Y$ 为 $R$ 上局部有限型概形之间的固有态射。则图
$$
\xymatrix{
Z_r(X_K) \ar[rr]_{sp_{X/R}} \ar[d] & & Z_r(X_\kappa) \ar[d] \\
Z_r(Y_K) \ar[rr]^{sp_{Y/R}} & & Z_r(Y_\kappa)
}
$$
可交换，其中 $r \geq 0$，竖直箭头由固有前推给出。
\end{lemma}

\begin{proof}
设 $Z \subset X$ 为支配 $R$ 的整闭子概形。由 $sp_{X/R}$ 的构造，有
$sp_{X/R}([Z_K]) = [Z_\kappa]_r$，而这刻画了特化映射。置
$Z' = f(Z) \subset Y$。则 $Z'$ 是 $Y$ 中支配 $R$ 的整闭子概形。因此
$sp_{Y/R}([Z'_K]) = [Z'_\kappa]_r$。

\medskip\noindent
可把 $[Z]$ 看成 $Z_{r + 1}(X)$ 的元素。按定义，若
$\dim(Z') < r + 1$，则 $f_*[Z] = 0$；若 $Z \to Z'$ 是次数为 $d$ 的
一般有限态射，则 $f_*[Z] = d[Z']$。由于固有前推与沿
$Y_K \to Y$ 的平坦拉回可交换（Chow 同调，引理
\ref{chow-lemma-flat-pullback-proper-pushforward}），相应地有
$f_{K, *}[Z_K] = 0$ 或 $f_{K, *}[Z_K] = d[Z'_K]$。对下列交换图应用
Chow 同调，引理 \ref{chow-lemma-closed-in-X-gysin}：
$$
\xymatrix{
X_\kappa \ar[d] \ar[r]_i & X \ar[d] \\
Y_\kappa \ar[r]^j & Y
}
$$
由此得到 $f_{\kappa, *}[Z_\kappa]_r = 0$，或
$f_{\kappa, *}[Z_\kappa] = d[Z'_\kappa]_r$，因为显然有
$i^*[Z] = [Z_k]_r$ 及 $j^*[Z'] = [Z'_\kappa]_r$。综合以上结果即得结论。
\end{proof}









\section{纤维上的循环族}
\label{relative-cycles-section-cycles-fibres}

\noindent
设 $f : X \to S$ 为概形的局部有限型态射，$r \geq 0$ 为整数。
{\it $X/S$ 各纤维上的 $r$-循环族 $\alpha$} 是一个族
$$
\alpha = (\alpha_s)_{s \in S}
$$
它以概形 $S$ 的点 $s$ 为指标，其中 $\alpha_s \in Z_r(X_s)$ 是 $f$ 在
$s$ 处的概形论纤维 $X_s$ 上的 $r$-循环。对纤维上的 $r$-循环族，可以
施行多种构造。

\medskip\noindent
{\bf 基变换。} 设
$$
\xymatrix{
X' \ar[r] \ar[d] & X \ar[d]^f \\
S' \ar[r]^g & S
}
$$
为概形态射的笛卡尔方块，且 $f$ 局部有限型。设 $r \geq 0$ 为整数。
给定 $X/S$ 各纤维上的 $r$-循环族 $\alpha$，定义 $\alpha$ 的
{\it 基变换} $g^*\alpha$ 为循环族
$$
g^*\alpha = (\alpha'_{s'})_{s' \in S'}
$$
其中 $\alpha'_{s'} \in Z_r(X'_{s'})$ 是循环 $\alpha_s$ 的基变换，
$s = g(s')$；这里按第 \ref{relative-cycles-section-relative-fields} 节，使用概形论纤维的
恒等式 $X'_{s'} = X_s \times_{\Spec(\kappa(s))} \Spec(\kappa(s'))$。

\medskip\noindent
{\bf 限制。} 设 $f : X \to S$ 为概形的局部有限型态射，$r \geq 0$ 为整数。
设 $U \subset X$ 与 $V \subset S$ 为开子概形，并且 $f(U) \subset V$。
给定 $X/S$ 各纤维上的 $r$-循环族 $\alpha$，可以把 $\alpha$ 的
{\it 限制} $\alpha|_U$ 定义为 $U/V$ 各纤维上的 $r$-循环族
$$
\alpha|_U = (\alpha_s|_{U_s})_{s \in V}
$$
即限制到各概形论纤维所得的循环族。

\medskip\noindent
{\bf 平坦拉回。} 设 $X \to S$ 为概形的局部有限型态射，$r, e \geq 0$
为整数。设 $f : X' \to X$ 为相对维数 $e$ 的局部有限型平坦态射。
给定 $X/S$ 各纤维上的 $r$-循环族 $\alpha$，定义 $\alpha$ 的
{\it 平坦拉回} $f^*\alpha$ 为各纤维上的 $(r + e)$-循环族
$$
f^*\alpha = (f_s^*\alpha_s)_{s \in S}
$$
其中 $f_s^*\alpha_s \in Z_{r + e}(X'_s)$ 是沿概形论纤维间相对维数为
$e$ 的平坦态射 $f_s : X'_s \to X_s$，对 $Z_r(X_s)$ 中的循环
$\alpha_s$ 作平坦拉回所得。

\medskip\noindent
{\bf 固有前推。} 设
$$
\xymatrix{
X \ar[rr]_f \ar[rd] & & Y \ar[ld] \\
& S
}
$$
为概形态射的交换图，其中 $X$ 与 $Y$ 在 $S$ 上局部有限型，且 $f$ 固有。
设 $r \geq 0$ 为整数。给定 $X/S$ 各纤维上的 $r$-循环族 $\alpha$，
定义 $\alpha$ 的{\it 固有前推} $f_*\alpha$ 为 $Y/S$ 各纤维上的
$r$-循环族
$$
f_*\alpha = (f_{s, *}\alpha_s)_{s \in S}
$$
其中 $f_{s, *}\alpha_s \in Z_r(Y_s)$ 是沿概形论纤维间的固有态射
$f_s : X_s \to Y_s$，对 $Z_r(X_s)$ 中的循环 $\alpha_s$ 作固有前推所得。

\begin{lemma}
\label{relative-cycles-lemma-compatibilities}
上述运算满足如下相容性：
(1) 基变换具有函子性；
(2) 限制是基变换与平坦拉回（的一个特殊情形）的复合；
(3) 平坦拉回与基变换可交换；
(4) 平坦拉回具有函子性；
(5) 固有前推与基变换可交换；
(6) 固有前推具有函子性；
(7) 固有前推与平坦拉回可交换。
\end{lemma}

\begin{proof}
把 Chow 同调一章中证明的相应结果应用于所涉概形在 $S$ 上的各纤维，
即可直接得到这些相容性。我们略去精确陈述与详细证明，只列出若干出处。
第 (1) 点：Chow 同调，引理 \ref{chow-lemma-compose-base-change}。
第 (2) 点：显然。
第 (3) 点：Chow 同调，引理
\ref{chow-lemma-pullback-base-change-pullback}。
第 (4) 点：Chow 同调，引理 \ref{chow-lemma-compose-flat-pullback}。
第 (5) 点：Chow 同调，引理
\ref{chow-lemma-pullback-base-change-pushforward}。
第 (6) 点：Chow 同调，引理 \ref{chow-lemma-compose-pushforward}。
第 (7) 点：Chow 同调，引理
\ref{chow-lemma-flat-pullback-proper-pushforward}。
\end{proof}

\begin{example}
\label{relative-cycles-example-family-associated-module}
设 $f : X \to S$ 为概形的局部有限型态射，$r \geq 0$ 为整数。设
$\mathcal{F}$ 为有限型拟凝聚 $\mathcal{O}_X$-模。对 $s \in S$，以
$\mathcal{F}_s$ 表示 $\mathcal{F}$ 到 $X_s$ 的拉回。假设对所有
$s \in S$ 都有 $\dim(\text{Supp}(\mathcal{F}_s)) \leq r$。于是可将
$X/S$ 各纤维上的 $r$-循环族 $[\mathcal{F}/X/S]_r$ 与
$\mathcal{F}$ 相伴；其定义公式为
$$
[\mathcal{F}/X/S]_r = ([\mathcal{F}_s]_r)_{s \in S}
$$
其中 $[\mathcal{F}_s]_r$ 由 Chow 同调，定义
\ref{chow-definition-cycle-associated-to-coherent-sheaf} 给出。
\end{example}

\begin{lemma}
\label{relative-cycles-lemma-family-associated-module}
例 \ref{relative-cycles-example-family-associated-module} 中的构造与基变换、限制和平坦拉回
相容。
\end{lemma}

\begin{proof}
见 Chow 同调，引理 \ref{chow-lemma-pullback-coherent-base-change} 与
\ref{chow-lemma-pullback-coherent}。
\end{proof}

\begin{example}
\label{relative-cycles-example-family-associated-closed}
设 $f : X \to S$ 为概形的局部有限型态射，$r \geq 0$ 为整数。设
$Z \subset X$ 为闭子概形。对 $s \in S$，以 $Z_s$ 表示 $Z$ 在 $X_s$
中的逆像；等价地，它是 $Z$ 在 $s$ 处的概形论纤维，并视为 $X_s$ 的闭
子概形。假设对所有 $s \in S$ 都有 $\dim(Z_s) \leq r$。于是可将
$X/S$ 各纤维上的 $r$-循环族 $[Z/X/S]_r$ 与 $Z$ 相伴；其定义公式为
$$
[Z/X/S]_r = ([Z_s]_r)_{s \in S}
$$
其中 $[Z_s]_r$ 由 Chow 同调，定义
\ref{chow-definition-cycle-associated-to-closed-subscheme} 给出。
\end{example}

\begin{lemma}
\label{relative-cycles-lemma-family-associated-closed}
例 \ref{relative-cycles-example-family-associated-closed} 中的构造与基变换、限制和平坦拉回
相容。
\end{lemma}

\begin{proof}
取 $\mathcal{F} = (Z \to X)_*\mathcal{O}_Z$，这就是引理
\ref{relative-cycles-lemma-family-associated-module} 的特殊情形。见 Chow 同调，引理
\ref{chow-lemma-cycle-closed-coherent}。
\end{proof}

\begin{remark}[支撑]
\label{relative-cycles-remark-supports-family}
设 $f : X \to S$ 为概形的局部有限型态射，$r \geq 0$ 为整数。设
$\alpha$ 为 $X/S$ 各纤维上的 $r$-循环族。定义 $\alpha$ 的{\it 支撑}为
$$
\text{Supp}(\alpha) =
\bigcup\nolimits_{s \in S} \text{Supp}(\alpha_s) \subset X
$$
这里 $\text{Supp}(\alpha_s) \subset X_s$ 是循环 $\alpha_s$ 的支撑；
见 Chow 同调，定义 \ref{chow-definition-support-cycle}。支撑
$\text{Supp}(\alpha)$ 很少是 $X$ 的闭子集。
\end{remark}

\begin{lemma}
\label{relative-cycles-lemma-support-family}
按评注 \ref{relative-cycles-remark-supports-family} 取支撑与基变换、限制和平坦拉回相容。
\end{lemma}

\begin{proof}
略。
\end{proof}

\begin{lemma}
\label{relative-cycles-lemma-coequalizer-dim-r}
设 $f : X \to S$ 为概形的局部有限型态射，$r \geq 0$ 为整数。设
$g : S' \to S$ 为概形的满射。令 $S'' = S' \times_S S'$，并以
$f' : X' \to S'$ 和 $f'' : X'' \to S''$ 表示 $f$ 的基变换。设
$x \in X$ 满足 $\text{trdeg}_{\kappa(f(x))}(\kappa(x)) = r$。
\begin{enumerate}
\item 存在映到 $x$ 的 $x' \in X'$，使得
$\text{trdeg}_{\kappa(f'(x'))}(\kappa(x')) = r$。
\item 若 $x'_1, x'_2 \in X'$ 均满足 (1)，则存在 $x'' \in X''$，使得
$\text{trdeg}_{\kappa(f''(x''))}(\kappa(x'')) = r$ 且
$\text{pr}_i(x'') = x'_i$。
\end{enumerate}
\end{lemma}

\begin{proof}
(1) 即“态射”引理
\ref{morphisms-lemma-dimension-fibre-after-base-change}。设
$x'_1, x'_2$ 如 (2) 所述。由于 $X'' = X' \times_X X'$，存在
$x'' \in X''$ 同时映到 $x'_1$ 和 $x'_2$（例如见“下降”引理
\ref{descent-lemma-equiv-fibre-product}）。分别以 $s'' \in S''$、
$s'_i \in S'$ 和 $s \in S$ 表示 $x''$、$x'_i$ 和 $x$ 的像。
记 $k = \kappa(s)$，并令 $Z \subset X_k$ 为以 $x$ 为一般点的整闭子概形。
于是 $x'_i$ 是 $Z_{\kappa(s'_i)}$ 某个不可约分支的一般点。取包含
$x''$ 的不可约分支 $Z'' \subset Z_{\kappa(s'')}$，并以
$\xi'' \in Z''$ 表示其一般点。由于 $\xi'' \leadsto x''$，在两个投影下，
$\xi''$ 也必映到 $x'_i$。另一方面，由于它是 $Z$ 的基变换的一个
不可约分支的一般点，故
$\text{trdeg}_{\kappa(s'')}(\kappa(\xi'')) = r$。
\end{proof}

\begin{lemma}
\label{relative-cycles-lemma-descend-family}
设 $f : X \to S$ 为概形的局部有限型态射，$r \geq 0$ 为整数。设
$g : S' \to S$ 为概形态射，并令 $X' = S' \times_S X$。假设对每个
$s \in S$，都存在点 $s' \in S'$，满足 $g(s') = s$，并且域扩张
$\kappa(s')/\kappa(s)$ 可分。则：
\begin{enumerate}
\item 对 $X/S$ 各纤维上的两个 $r$-循环族 $\alpha_1$ 和 $\alpha_2$，若
$g^*\alpha_1 = g^*\alpha_2$，则 $\alpha_1 = \alpha_2$。
\item 给定 $X'/S'$ 各纤维上的 $r$-循环族 $\alpha'$，若作为
$(S' \times_S S') \times_S X / (S' \times_S S')$ 各纤维上的
$r$-循环族，有 $\text{pr}_1^*\alpha' = \text{pr}_2^*\alpha'$，则存在唯一的
$X/S$ 各纤维上的 $r$-循环族 $\alpha$，使得 $g^*\alpha = \alpha'$。
\end{enumerate}
\end{lemma}

\begin{proof}
(1) 由第 \ref{relative-cycles-section-relative-fields} 节所述基变换映射的单射性得出。
（只要 $S' \to S$ 满射，这一论证就成立。）

\medskip\noindent
设 $\alpha'$ 如 (2) 所述，并以
$\alpha'' = \text{pr}_1^*\alpha' = \text{pr}_2^*\alpha'$
表示其公共值。

\medskip\noindent
令 $(X/S)^{(r)}$ 为所有满足
$\text{trdeg}_{\kappa(f(x))}(\kappa(x)) = r$ 的 $x \in X$ 所成之集，
并类似定义 $(X'/S')^{(r)}$ 和 $(X''/S'')^{(r)}$。取系数后，可以把
$\alpha'$ 和 $\alpha''$ 看成函数
$\alpha' : (X'/S')^{(r)} \to \mathbf{Z}$ 和
$\alpha'' : (X''/S'')^{(r)} \to \mathbf{Z}$。给定函数
$$
\varphi : (X/S)^{(r)} \to \mathbf{Z}
$$
仿照基变换运算，定义
$g^*\varphi : (X'/S')^{(r)} \to \mathbf{Z}$。具体地，设
$x' \in (X'/S')^{(r)}$ 映到 $x \in X$、$s' \in S'$ 和 $s \in Z$。
分别以 $Z' \subset X'_{s'}$ 和 $Z \subset X_s$ 表示以 $x'$ 和 $x$ 为
一般点的整闭子概形。注意 $\dim(Z') = r$。若 $\dim(Z) < r$，则令
$(g^*\varphi)(x') = 0$。若 $\dim(Z) = r$，则 $Z'$ 是 $Z_{s'}$ 的一个
不可约分支，因而具有重数 $m_{Z', Z_{s'}}$；将其记作 $m(x', g)$。
定义
$$
(g^*\varphi)(x') = m(x', g) \varphi(x)
$$
注意，系数 $m(x', g)$ 总是正整数（例如见引理
\ref{relative-cycles-lemma-multiplicities-field-extension}）。类似地，有基变换映射
$$
\text{pr}_1^*, \text{pr}_2^* :
\text{Map}((X'/S')^{(r)}, \mathbf{Z})
\longrightarrow
\text{Map}((X''/S'')^{(r)}, \mathbf{Z})
$$
由基变换的结合性，
$\text{pr}_1^* \circ g^* = \text{pr}_2^* \circ g^*$（略去一个小细节）。
明确地说，就集合映射而言，该等式意味着对
$x'' \in (X''/S'')^{(r)}$ 有
$$
m(x'', \text{pr}_1) m(\text{pr}_1(x''), g) =
m(x'', \text{pr}_2) m(\text{pr}_2(x''), g)
$$
只要 $\text{pr}_1(x'')$ 和 $\text{pr}_2(x'')$ 属于
$(X''/S'')^{(r)}$。由引理 \ref{relative-cycles-lemma-coequalizer-dim-r}、上述等式和一个
初等论证\footnote{给定 $x \in (X/S)^{(r)}$，选取映到 $x$ 的
$x' \in (X'/S')^{(r)}$，并令
$\alpha(x) = \alpha'(x')/m(x', g)$。由该公式和引理，这一取值良定。}，
存在唯一映射
$$
\alpha : (X/S)^{(r)} \to \mathbf{Q}
$$
使得 $g^*\alpha = \alpha'$。为完成证明，只需说明 $\alpha$ 取整数值
（略去一个小细节：还须说明 $\alpha$ 在每条纤维上确定局部有限和；这由
$\alpha'$ 的相应性质得出）。给定像为 $s \in S$ 的任意
$x \in (X/S)^{(r)}$，可以选取点 $s' \in S'$，使得
$\kappa(s')/\kappa(s)$ 可分。于是可选取 $x' \in (X'/S')^{(r)}$，
使其映到 $s$ 和 $x$；此时 $Z_{s'}$ 约化，故 $m(x', g) = 1$。
因此 $\alpha(x) = \alpha'(x')$ 是整数。
\end{proof}

\begin{lemma}
\label{relative-cycles-lemma-pullback-universally-bijective}
设 $g : S' \to S$ 为诱导剩余域同构的概形双射态射。设
$f : X \to S$ 局部有限型，并令 $X' = S' \times_S X$。设 $r \geq 0$。
则沿 $g$ 的基变换在 $X/S$ 各纤维上的 $r$-循环族之群与
$X'/S'$ 各纤维上的 $r$-循环族之群之间确定一个双射。
\end{lemma}

\begin{proof}
略。
\end{proof}








\section{相对循环}
\label{relative-cycles-section-families-specialization}

\noindent
我们将使用如下定义；它与 \cite{SV} 中定义的比较见第
\ref{relative-cycles-section-compare} 节。

\begin{definition}
\label{relative-cycles-definition-relative-cycles}
设 $S$ 为局部 Noether 概形，$f : X \to S$ 为概形的局部有限型态射，
$r \geq 0$ 为整数。$X/S$ 上的一个{\it 相对 $r$-循环}，是指
$X/S$ 各纤维上的一个 $r$-循环族 $\alpha$，使得对每个态射
$g : S' \to S$，只要 $S'$ 是离散赋值环的谱，就有
$$
sp_{X'/S'}(\alpha_\eta) = \alpha_0
$$
其中 $sp_{X'/S'}$ 如第 \ref{relative-cycles-section-specialization} 节所定义，而
$\alpha_\eta$（相应地，$\alpha_0$）是 $\alpha$ 的基变换
$g^*\alpha$ 在 $S'$ 的一般点（相应地，闭点）处的值。$X/S$ 上所有
相对 $r$-循环所成的群记作 $z(X/S, r)$。
\end{definition}

\begin{lemma}
\label{relative-cycles-lemma-relative-cycle-functoriality}
设 $\alpha$ 是定义 \ref{relative-cycles-definition-relative-cycles} 所述的 $X/S$ 上相对
$r$-循环。则 $\alpha$ 的任意限制、基变换、平坦拉回或固有前推仍为相对
$r$-循环。
\end{lemma}

\begin{proof}
平坦拉回的情形使用引理 \ref{relative-cycles-lemma-specialization-flat-pullback}。
限制是平坦拉回的特殊情形。基变换的情形使用基变换的传递性。
固有前推的情形使用引理 \ref{relative-cycles-lemma-specialization-proper-pushforward}。
\end{proof}

\begin{lemma}
\label{relative-cycles-lemma-relative-cycles-h-descent}
设 $f : X \to S$ 为概形态射。假设 $S$ 局部 Noether，且 $f$ 局部有限型。
设 $r \geq 0$ 为整数，$\alpha$ 为 $X/S$ 各纤维上的 $r$-循环族。设
$\{g_i : S_i \to S\}$ 为一个 h-覆盖（“平坦性续篇”定义
\ref{flat-definition-h-covering}）。则 $\alpha$ 是相对 $r$-循环，当且仅当
每个基变换 $g_i^*\alpha$ 都是相对 $r$-循环。
\end{lemma}

\begin{proof}
若 $\alpha$ 是相对 $r$-循环，则由引理
\ref{relative-cycles-lemma-relative-cycle-functoriality}，每个基变换 $g_i^*\alpha$ 都是相对
$r$-循环。反之，假设每个 $g_i^*\alpha$ 都是相对 $r$-循环。设
$g : S' \to S$ 为态射，其中 $S'$ 是离散赋值环的谱。把 $S$ 替换为
$S'$、把 $X$ 替换为 $X' = X \times_S S'$、把 $\alpha$ 替换为
$\alpha' = g^*\alpha$，并利用 h-覆盖的基变换仍为 h-覆盖（“平坦性续篇”
引理 \ref{flat-lemma-h}），便归约到下一段所研究的问题。

\medskip\noindent
假设 $S$ 是离散赋值环的谱，其闭点为 $0$、一般点为 $\eta$。需要证明
$sp_{X/S}(\alpha_\eta) = \alpha_0$。由于 h-覆盖按定义是 V-覆盖，存在
$i$ 以及 $S_i$ 中点的特化 $s' \leadsto s$，满足
$g_i(s') = \eta$ 和 $g_i(s) = 0$；见“拓扑”引理
\ref{topologies-lemma-refine-qcqs-V}。由“性质”引理
\ref{properties-lemma-locally-Noetherian-specialization-dvr}，存在从离散
赋值环的谱 $S'$ 出发的态射 $h : S' \to S_i$，它把一般点 $\eta'$ 映到
$s'$，把闭点 $0'$ 映到 $s$。记 $\alpha' = h^*g_i^*\alpha$。由假设，
$sp_{X'/S'}(\alpha'_{\eta'}) = \alpha'_{0'}$。由于
$g = g_i \circ h : S' \to S$ 是由离散赋值环扩张诱导的概形态射，
$sp_{X/S}$ 与 $sp_{X'/S'}$ 和纤维上的基变换映射相容；见引理
\ref{relative-cycles-lemma-specialization-extension}。第 \ref{relative-cycles-section-relative-fields} 节
已经说明基变换映射 $Z_r(X_0) \to Z_r(X'_{0'})$ 是单射，故
$sp_{X/S}(\alpha_\eta) = \alpha_0$。
\end{proof}

\begin{lemma}
\label{relative-cycles-lemma-families-specialization-fppf-descent}
设 $f : X \to S$ 为概形态射。假设 $S$ 局部 Noether，且 $f$ 局部有限型。
设 $r, e \geq 0$ 为整数，$\alpha$ 为 $X/S$ 各纤维上的 $r$-循环族。
设 $\{f_i : X_i \to X\}$ 为一族联合满射的局部有限型平坦态射，且其相对
维数为 $e$。则 $\alpha$ 是相对 $r$-循环，当且仅当每个平坦拉回
$f_i^*\alpha$ 都是相对 $r$-循环。
\end{lemma}

\begin{proof}
若 $\alpha$ 是相对 $r$-循环，则由引理
\ref{relative-cycles-lemma-relative-cycle-functoriality}，每个拉回 $f_i^*\alpha$ 都是相对
$r$-循环。反之，假设每个 $f_i^*\alpha$ 都是相对 $r$-循环。设
$g : S' \to S$ 为态射，其中 $S'$ 是离散赋值环的谱。把 $S$ 替换为
$S'$、把 $X$ 替换为 $X' = X \times_S S'$、把 $\alpha$ 替换为
$\alpha' = g^*\alpha$，便归约到下一段所研究的问题。

\medskip\noindent
假设 $S$ 是离散赋值环的谱，其闭点为 $0$、一般点为 $\eta$。需要证明
$sp_{X/S}(\alpha_\eta) = \alpha_0$。以
$f_{i, 0} : X_{i, 0} \to X_0$ 表示 $f_i$ 到 $S$ 的闭点上的基变换；
$f_{i, \eta}$ 的记号类似。注意
$$
f_{i, 0}^*sp_{X/S}(\alpha_\eta) =
sp_{X_i/S}(f_{i, \eta}^*\alpha_\eta) = f_{i, 0}^*\alpha_0
$$
其中第一个等式由引理 \ref{relative-cycles-lemma-specialization-flat-pullback} 得出，
第二个等式由假设得出。由于映射族
$f_{i, 0}^* : Z_r(X_0) \to Z_r(X_{i, 0})$ 联合单射（这是因为
$f_{i, 0}$ 联合满射），结论随即成立。
\end{proof}

\begin{lemma}
\label{relative-cycles-lemma-check-after-closed}
设 $S$ 为局部 Noether 概形，$i : X \to Y$ 为 $S$ 上局部有限型概形之间的
闭浸入。设 $r \geq 0$，并令 $\alpha$ 为 $X/S$ 各纤维上的 $r$-循环族。
则 $\alpha$ 是 $X/S$ 上的相对 $r$-循环，当且仅当 $i_*\alpha$ 是
$Y/S$ 上的相对 $r$-循环。
\end{lemma}

\begin{proof}
由于基变换与 $i_*$ 可交换（引理 \ref{relative-cycles-lemma-compatibilities}），只需证明：
若 $S$ 是离散赋值环的谱，其一般点为 $\eta$、闭点为 $0$，则
$sp_{X/S}(\alpha_\eta) = \alpha_0$ 当且仅当
$sp_{Y/S}(i_{\eta, *}\alpha_\eta) = i_{0, *}\alpha_0$。这是因为
$i_{0, *} : Z_r(X_0) \to Z_r(Y_0)$ 是单射，而且由引理
\ref{relative-cycles-lemma-specialization-proper-pushforward}，有
$i_{0, *}sp_{X/S}(\alpha_\eta) =
sp_{Y/S}(i_{\eta, *}\alpha_\eta)$。
\end{proof}

\noindent
下面的引理将在引理 \ref{relative-cycles-lemma-relative-cycles-equal} 中得到加强。

\begin{lemma}
\label{relative-cycles-lemma-uniqueness-extension}
设 $f : X \to S$ 为概形态射。假设 $S$ 局部 Noether，且 $f$ 局部有限型。
设 $r \geq 0$，并令 $\alpha$ 和 $\beta$ 为 $X/S$ 上的相对 $r$-循环。
下列条件等价：
\begin{enumerate}
\item $\alpha = \beta$；
\item 对 $S$ 的任一不可约分支的一般点 $\eta \in S$，都有
$\alpha_\eta = \beta_\eta$。
\end{enumerate}
\end{lemma}

\begin{proof}
(1) $\Rightarrow$ (2) 显然。假设 (2)。对每个 $s \in S$，都可以找到
(2) 所述的一个 $\eta$，使其特化到 $s$。由“性质”引理
\ref{properties-lemma-locally-Noetherian-specialization-dvr}，存在从离散
赋值环的谱 $S'$ 出发的态射 $g : S' \to S$，它把一般点 $\eta'$ 映到
$\eta$，把闭点 $0$ 映到 $s$。于是 $\alpha_s$ 和 $\beta_s$ 是
$Z_r(X_s)$ 中的元素，并基变换为 $Z_r(X_{0'})$ 中同一个元素，即
$sp_{X_{S'}/S'}(\alpha_{\eta'})$，其中 $\alpha_{\eta'}$ 是
$\alpha_\eta$ 的基变换。第 \ref{relative-cycles-section-relative-fields} 节已经说明
基变换映射 $Z_r(X_s) \to Z_r(X_{0'})$ 是单射，故
$\alpha_s = \beta_s$。
\end{proof}

\begin{lemma}
\label{relative-cycles-lemma-family-associated-module-specialization}
在例 \ref{relative-cycles-example-family-associated-module} 的情形中，假设 $S$ 局部
Noether，且 $\mathcal{F}$ 在维数 $\geq r$ 处在 $S$ 上平坦
（“平坦性续篇”定义 \ref{flat-definition-flat-dimension-n}）。则
$[\mathcal{F}/X/S]_r$ 是 $X/S$ 上的相对 $r$-循环。
\end{lemma}

\begin{proof}
由“平坦性续篇”引理 \ref{flat-lemma-pre-flat-dimension-n}，关于
$\mathcal{F}$ 的假设在任意基变换下保持成立。此外，由引理
\ref{relative-cycles-lemma-family-associated-module}，$[\mathcal{F}/X/S]_r$ 的构造与任意
基变换相容。由于与特化相容这一条件是在基变换到离散赋值环的谱之后
检验的，问题归约到 $S$ 是赋值环的谱的情形。此时集合
$U = \{x \in X \mid \mathcal{F}\text{ 在点 }x\text{ 上平坦，底为 }S\}$
由“平坦性续篇”引理 \ref{flat-lemma-finite-type-flat} 是 $X$ 中的开集。
根据假设，$X$ 中 $U$ 的补集在 $S$ 上各纤维的维数 $< r$，故沿浸入
$U \subset X$ 的限制在任意基变换后都诱导纤维上 $r$-循环群的同构，
并与特化映射及由 $\mathcal{F}$ 所伴随相对循环的构造相容。因此，只需证明
$[\mathcal{F}|_U / U /S]_r$ 与特化相容。由于 $\mathcal{F}|_U$ 在 $S$
上平坦，这由引理 \ref{relative-cycles-lemma-specialization-module} 和定义得出。
\end{proof}

\begin{lemma}
\label{relative-cycles-lemma-family-associated-closed-specialization}
在例 \ref{relative-cycles-example-family-associated-closed} 的情形中，假设 $S$ 局部
Noether，且 $Z$ 在维数 $\geq r$ 处在 $S$ 上平坦。则
$[Z/X/S]_r$ 是 $X/S$ 上的相对 $r$-循环。
\end{lemma}

\begin{proof}
该假设意味着 $\mathcal{O}_Z$ 在维数 $\geq r$ 处在 $S$ 上平坦。于是对
$\mathcal{F} = (Z \to X)_*\mathcal{O}_Z$ 应用引理
\ref{relative-cycles-lemma-family-associated-module-specialization} 即得结论。
\end{proof}

\noindent
设 $S$ 为局部 Noether 概形，$f : X \to S$ 为有限型态射，并设
$r \geq 0$。以 $Hilb(X/S, r)$ 表示所有闭子概形 $Z \subset X$ 所成的
集合，其中 $Z \to S$ 平坦且相对维数 $\leq r$。由引理
\ref{relative-cycles-lemma-family-associated-closed-specialization}，对每个
$Z \in Hilb(X/S, r)$，都有元素 $[Z/X/S]_r \in z(X/S, r)$。因此得到群同态
\begin{equation}
\label{relative-cycles-equation-cycle-classes}
\text{由下列集合生成的自由 Abel 群： }Hilb(X/S, r) \longrightarrow z(X/S, r)
\end{equation}
它把 $\sum n_i[Z_i]$ 映到 $\sum n_i[Z_i/X/S]_r$。相对 $r$-循环的一个
关键特征是：在适当拓扑下，它们局部地（在 $X$ 和 $S$ 上）属于该映射的像。

\begin{lemma}
\label{relative-cycles-lemma-get-cycles}
设 $f : X \to S$ 为概形的有限型态射，且 $S$ Noether。设 $r \geq 0$，
$\alpha$ 为 $X/S$ 上的相对 $r$-循环。则存在固有完全分解态射
（“态射续篇”定义 \ref{more-morphisms-definition-cd-morphism}）
$g : S' \to S$，使得 $g^*\alpha$ 属于 (\ref{relative-cycles-equation-cycle-classes}) 的像。
\end{lemma}

\begin{proof}
由 Noether 归纳法，可以假设：对任意不是同构的闭浸入
$g : S' \to S$，结论对 $\alpha$ 的拉回成立。

\medskip\noindent
令 $S_1 \subset S$ 为一个不可约分支（视作整闭子概形）。令
$S_2 \subset S$ 为 $S'$ 的补集之闭包（视作约化闭子概形）。若
$S_2 \not = \emptyset$，则结论对沿 $S_1 \to S$ 和 $S_2 \to S$ 的
$\alpha$ 拉回成立。若 $g_1 : S'_1 \to S_1$ 和 $g_2 : S'_2 \to S_2$
是相应的固有完全分解态射，则
$S' = S'_1 \amalg S'_2 \to S$ 是固有完全分解态射，且结论对 $S$
成立\footnote{确实，$S'_1 \times_S X$ 中任一在 $S'_1$ 上平坦且相对维数
$\leq r$ 的闭子概形，都可以视作 $S' \times_S X$ 中一个在 $S'$ 上平坦
且相对维数 $\leq r$ 的闭子概形。}。因此可以假设 $S' \to S$ 是双射，
从而归约到下一段所述情形。

\medskip\noindent
假设 $S$ 整。令 $\eta \in S$ 为一般点，并令 $K = \kappa(\eta)$ 为
$S$ 的函数域。于是 $\alpha_\eta$ 是 $X_K$ 上的 $r$-循环。写成
$\alpha_\eta = \sum n_i[Y_i]$。取 $Y_i$ 的闭包，得到整闭子概形
$Z_i \subset X$，其到 $\eta$ 的基变换为 $Y_i$。由一般平坦性（例如
“态射”命题 \ref{morphisms-proposition-generic-flatness}），对每个 $i$，
$Z_i$ 在 $S$ 的某个非空开集 $U$ 上平坦。应用“平坦性续篇”引理
\ref{flat-lemma-flat-after-blowing-up}，可以找到 $U$-容许的吹起
$g : S' \to S$，使得 $Z_i$ 的严格变换 $Z'_i \subset X_{S'}$ 在 $S'$
上平坦。于是 $\beta = \sum n_i[Z'_i/X_{S'}/S']_r$ 属于
(\ref{relative-cycles-equation-cycle-classes}) 的像；由引理
\ref{relative-cycles-lemma-uniqueness-extension}，$\beta = g^*\alpha$。

\medskip\noindent
但这尚未完成证明，因为 $S' \to S$ 未必完全分解。修正这一点很容易：
以 $T \subset S$ 表示 $U$ 的补集（视作闭子概形）。由 Noether 归纳法，
可以找到固有完全分解态射 $T' \to T$，使得
$(T' \to S)^*\alpha$ 属于 (\ref{relative-cycles-equation-cycle-classes}) 的像。于是
$S' \amalg T' \to S$ 满足要求。
\end{proof}

\begin{lemma}
\label{relative-cycles-lemma-get-cycles-dvr}
设 $f : X \to S$ 为概形的有限型态射，其中 $S$ 是离散赋值环的谱。
设 $r \geq 0$。则 (\ref{relative-cycles-equation-cycle-classes}) 是满射。
\end{lemma}

\begin{proof}
这当然由引理 \ref{relative-cycles-lemma-get-cycles} 得出，也可以如下直接看出。设
$\alpha$ 是 $X/S$ 上的相对 $r$-循环。写
$\alpha_\eta = \sum n_i[Z_i]$（该和有限）。如第
\ref{relative-cycles-section-specialization} 节，以 $\overline{Z}_i \subset X$ 表示
$Z_i$ 的闭包。于是 $\alpha = \sum n_i[\overline{Z}_i/X/S]$。
\end{proof}

\begin{lemma}
\label{relative-cycles-lemma-relative-cycle-smooth}
设 $f : X \to S$ 为概形态射，$r \geq 0$。假设 $S$ 局部 Noether，且
$f$ 光滑、相对维数为 $r$。设 $\alpha \in z(X/S, r)$。则 $\alpha$ 的支撑
在 $X$ 中既开又闭（更精确的结论见证明）。
\end{lemma}

\begin{proof}
设 $x \in X$ 的像为 $s \in S$。由于 $f$ 光滑，$X_s$ 中存在唯一一个包含
$x$ 的不可约分支 $Z(x)$。于是 $\dim(Z(x)) = r$。以 $n_x$ 表示循环
$\alpha_s$ 中 $Z(x)$ 的系数。我们将证明函数 $x \mapsto n_x$ 在 $X$ 上
局部常值。

\medskip\noindent
设 $g : S' \to S$ 为局部 Noether 概形之间的态射。以 $X'$ 表示 $X$ 的
基变换，以 $\alpha' = g^*\alpha$ 表示 $\alpha$ 的基变换。设
$x' \in X'$ 映到 $s' \in S'$、$x \in X$ 和 $s \in S$。我们断言
$n_{x'} = n_x$。事实上，由于 $Z(x)$ 在 $\kappa(s)$ 上光滑，概形
$Z(x) \times_{\Spec(\kappa(s))} \Spec(\kappa(s'))$
是约化的。由于 $Z(x')$ 是该概形的一个不可约分支，根据第
\ref{relative-cycles-section-relative-fields} 节的基变换定义，$\alpha'_{s'}$ 中
$Z(x')$ 的系数 $n_{x'}$ 等于 $\alpha_s$ 中 $Z(x)$ 的系数 $n_x$，
从而证明了该断言。

\medskip\noindent
由于 $X$ 局部 Noether，为证明 $x \mapsto n_x$ 局部常值，只需证明：
若 $x' \leadsto x$ 是 $X$ 中的一个特化，则 $n_{x'} = n_x$。选取态射
$S' \to X$，其中 $S'$ 是离散赋值环的谱，它把一般点 $\eta$ 映到 $x'$，
把闭点 $0$ 映到 $x$；见“性质”引理
\ref{properties-lemma-locally-Noetherian-specialization-dvr}。于是 $f$ 沿
$S' \to S$ 的基变换 $X' \to S'$ 有一个截面 $\sigma : S' \to X'$，使得
$\sigma(\eta) \leadsto \sigma(0)$ 是 $X'$ 中点的特化，并映到 $X$ 中的
$x' \leadsto x$。因此归约到下一段的断言。

\medskip\noindent
设 $S$ 是离散赋值环的谱，其一般点为 $\eta$、闭点为 $0$，并且给定截面
$\sigma : S  \to X$。断言：$n_{\sigma(\eta)} = n_{\sigma(0)}$。由
“态射续篇”第 \ref{more-morphisms-section-connected-components} 节中的
讨论，特别是引理
\ref{more-morphisms-lemma-connected-along-section-open}，把 $X$ 替换为一个
开子概形后，可以假设 $X \to S$ 的纤维连通。由于这些纤维光滑，它们
不可约。于是 $\alpha_\eta = n[X_\eta]$，其中 $n = n_{\sigma(\eta)}$；关系
$sp_{X/S}(\alpha_\eta) = \alpha_0$ 蕴含 $\alpha_0 = n[X_0]$，即所需的
$n_{\sigma(0)} = n$。
\end{proof}

\begin{lemma}
\label{relative-cycles-lemma-relative-cycles-equal}
设 $f : X \to S$ 为概形态射。假设 $S$ 局部 Noether，且 $f$ 局部有限型。
设 $r \geq 0$ 且 $\alpha, \beta \in z(X/S, r)$。则集合
$E = \{s \in S : \alpha_s = \beta_s\}$ 在 $S$ 中闭。
\end{lemma}

\begin{proof}
问题在 $S$ 上是局部的，故可假设 $S$ 仿射。取仿射开覆盖
$X = \bigcup U_i$，并令
$E_i = \{s \in S : \alpha_s|_{U_{i, s}} = \beta_s|_{U_{i, s}}\}$。
于是 $E = \bigcap E_i$。因此只需对 $U_i \to S$ 以及 $\alpha$、$\beta$
到 $U_i$ 的限制证明该引理。这就归约到下一段所述情形。

\medskip\noindent
假设 $X$ 和 $S$ 拟紧。令 $\gamma = \alpha - \beta$，则
$E = \{s \in S : \gamma_s = 0\}$。由引理
\ref{relative-cycles-lemma-family-associated-closed-specialization}，存在有限族联合满射的
固有态射 $\{g_i : S_i \to S\}$，使得 $g_i^*\gamma$ 属于
(\ref{relative-cycles-equation-cycle-classes}) 的像。注意，$E_i = g_i^{-1}(E)$ 是所有满足
$(g_i^*\gamma)_t = 0$ 的点 $t \in S_i$ 所成之集。若对所有 $i$，$E_i$ 都闭，
则 $E = \bigcup g_i(E_i)$ 也闭。这就归约到下一段所述情形。

\medskip\noindent
假设 $X$ 和 $S$ 拟紧，且
$\gamma = \sum n_i[Z_i/X/S]_r$，其中只有有限个闭子概形 $Z_i \subset X$，
它们在 $S$ 上平坦且相对维数 $\leq r$。令
$X' = \bigcup Z_i$（概形论并）。于是 $i : X' \to X$ 是闭浸入，且
$X'$ 在 $S$ 上的相对维数 $\leq r$。此外，$\gamma = i_*\gamma'$，其中
$\gamma' = \sum n_i[Z_i/X'/S]_r$。显然
$E = E' = \{s \in S : \gamma'_s = 0\}$，故归约到下一段所述情形。

\medskip\noindent
假设 $X$ 在 $S$ 上的相对维数 $\leq r$。取 $s \in S$ 且
$s \not \in E$。我们将证明存在 $s$ 的开邻域 $V \subset S$，使得
$E \cap V$ 为空。假设 $s \not \in E$ 意味着存在维数为 $r$ 的整闭子概形
$Z \subset X_s$，使得 $\gamma_s$ 中 $[Z]$ 的系数 $n$ 非零。令
$x \in Z$ 为一般点。由于 $\dim(Z) = r$，$x$ 是 $X_s$ 的一个不可约分支
（即 $Z$）的一般点。因此，把 $X$ 替换为 $x$ 的一个开邻域后，可以假设
$Z$ 是 $X_s$ 的唯一不可约分支。特别地，$\gamma_s = n[Z]$。

\medskip\noindent
现在应用“态射续篇”引理
\ref{more-morphisms-lemma-local-structure-finite-type}，得到图
$$
\xymatrix{
X \ar[dd] & X' \ar[l]^g \ar[d]^\pi & x \ar@{|->}[dd] &
x' \ar@{|->}[l]  \ar@{|->}[d] \\
& Y \ar[d]^h & & y \ar@{|->}[d] \\
S \ar@{=}[r] & S & s & s \ar@{=}[l]
}
$$
并具有该引理列出的全部性质。令 $\gamma' = g^*\gamma$ 为平坦拉回。
注意 $E \subset E' = \{s \in S: \gamma'_s = 0\}$，并且
$s \not \in E'$，因为 $\gamma'_s$ 中 $Z'$ 的系数非零；这里
$Z' \subset X'_s$ 是 $x'$ 的闭包。类似地，令
$\gamma'' = \pi_*\gamma'$。于是
$E' \subset E'' = \{s \in S: \gamma''_s = 0\}$，且
$s \not \in E''$，因为 $\gamma''_s$ 中 $Z''$ 的系数非零；这里
$Z'' \subset Y_s$ 是 $y$ 的闭包。由引理
\ref{relative-cycles-lemma-relative-cycle-smooth} 和 $Y \to S$ 的开性，$s$ 的某个开邻域
与 $E''$ 不交，证明完成。
\end{proof}

\begin{lemma}
\label{relative-cycles-lemma-descend-through-limit}
设 $S = \lim_{i \in I} S_i$ 是一个 Noether 概形有向逆系统的极限，且转移
态射均为仿射态射。设 $0 \in I$，并令 $X_0 \to S_0$ 为概形的有限型态射。
对 $i \geq 0$，令 $X_i = S_i \times_{S_0} X_0$，并令
$X = S \times_{S_0} X_0$。若 $S$ 也 Noether，则
$$
z(X/S, r) = \colim_{i \geq 0} z(X_i/S_i, r)
$$
其中转移映射由相对 $r$-循环的基变换给出。
\end{lemma}

\begin{proof}
假设 $i \geq 0$ 且 $\alpha_i, \beta_i \in z(X_i/S_i, r)$ 映到
$z(X/S, r)$ 中同一个元素。则 $S \to S_i$ 的像包含在引理
\ref{relative-cycles-lemma-relative-cycles-equal} 的闭子集 $E \subset S_i$ 中。因此，对某个
$j \geq i$，态射 $S_j \to S_i$ 的像包含在 $E$ 中；见“极限”引理
\ref{limits-lemma-limit-contained-in-constructible}。所以 $\alpha_i$ 和
$\beta_i$ 到 $S_j$ 的基变换相等。故该映射为单射。

\medskip\noindent
设 $\alpha \in z(X/S, r)$。应用引理 \ref{relative-cycles-lemma-get-cycles}，得到固有完全
分解态射 $g : S' \to S$，使得 $g^*\alpha$ 属于
(\ref{relative-cycles-equation-cycle-classes}) 的像。令 $X' = S' \times_S X$。可以写成
$g^*\alpha = \sum n_a [Z_a/X'/S']_r$，其中 $Z_a \subset X'$ 是在 $S'$
上平坦且相对维数 $\leq r$ 的闭子概形。

\medskip\noindent
现在运用“极限”第 \ref{limits-section-descending-relative} 节的工具。
可以找到 $i \geq 0$，使得存在：
\begin{enumerate}
\item 固有完全分解态射 $g_i : S'_i \to S_i$，它到 $S$ 的基变换为
$g : S' \to S$；
\item 令 $X'_i = S'_i \times_{S_i} X_i$，则有闭子概形
$Z_{ai} \subset X'_i$，它们在 $S'_i$ 上平坦、相对维数 $\leq r$，且到
$S'$ 的基变换为 $Z_a$。
\end{enumerate}
这里使用了“极限”引理
\ref{limits-lemma-descend-finite-presentation},
\ref{limits-lemma-descend-closed-immersion-finite-presentation},
\ref{limits-lemma-descend-flat-finite-presentation},
\ref{limits-lemma-eventually-proper} 和
\ref{limits-lemma-limit-dimension}
以及“态射续篇”引理 \ref{more-morphisms-lemma-descend-cd}。考虑
$\alpha'_i = \sum n_a [Z_{ai}/X'_i/S'_i]_r \in z(X'_i/S'_i, r)$.
依构造，$\alpha'_i$ 在 $z(X'/S', r)$ 中的像等于基变换 $g^*\alpha$。

\medskip\noindent
令 $S''_i = S'_i \times_{S_i} S'_i$、
$X''_i = S''_i \times_{S_i} X_i$，并令 $S'' = S' \times_S S'$、
$X'' = S'' \times_S X$。以 $\text{pr}_1, \text{pr}_2 : S'' \to S'$ 和
$\text{pr}_1, \text{pr}_2 : S''_i \to S'_i$ 表示投影。两个基变换
$\text{pr}_1^*\alpha'_i$ 和 $\text{pr}_1^*\alpha'_i$ 映到
$z(X''/S'', r)$ 中同一个元素，因为
$\text{pr}_1^*g^*\alpha = \text{pr}_1^*g^*\alpha$。因此增大 $i$ 后，
由证明的第一段，可以假设
$\text{pr}_1^*\alpha'_i = \text{pr}_1^*\alpha'_i$。由引理
\ref{relative-cycles-lemma-descend-family}，得到 $X_i/S_i$ 各纤维上唯一的 $r$-循环族
$\alpha_i$，使得 $g_i^*\alpha_i = \alpha'_i$（这里用到
$S'_i \to S_i$ 完全分解）。由引理
\ref{relative-cycles-lemma-relative-cycles-h-descent}，$\alpha_i \in z(X_i/S_i, r)$。
引理 \ref{relative-cycles-lemma-descend-family} 中的唯一性说明 $\alpha_i$ 在
$z(X/S, r)$ 中的像为 $\alpha$，证明完成。
\end{proof}

\begin{lemma}
\label{relative-cycles-lemma-thickening}
设 $S$ 为局部 Noether 概形，$i : X \to X'$ 为 $S$ 上局部有限型概形之间的
加厚。设 $r \geq 0$。则 $i_* : z(X/S, r) \to z(X'/S, r)$ 是双射。
\end{lemma}

\begin{proof}
由于 $i_s : X_s \to X'_s$ 是加厚，显然 $i_*$ 在 $X/S$ 各纤维上的
$r$-循环族与 $X'/S$ 各纤维上的 $r$-循环族之间诱导双射。此外，对
$X/S$ 各纤维上的 $r$-循环族 $\alpha$，由引理
\ref{relative-cycles-lemma-check-after-closed}，
$\alpha \in z(X/S, r) \Leftrightarrow i_*\alpha \in z(X'/S, r)$。
引理得证。
\end{proof}

\begin{lemma}
\label{relative-cycles-lemma-extend-to-larger}
设 $S$ 为局部 Noether 概形，$X$ 为 $S$ 上局部有限型概形。设
$r \geq 0$。设 $U \subset X$ 为开集，使得 $X \setminus U$ 在 $S$ 上的
相对维数 $< r$；换言之，对所有 $s \in S$，都有
$\dim(X_s \setminus U_s) < r$。则限制确定双射
$z(X/S, r) \to z(U/S, r)$。
\end{lemma}

\begin{proof}
由维数假设，$Z_r(X_s) \to Z_r(U_s)$ 是双射，所以限制在 $X/S$ 各纤维上的
$r$-循环族与 $U/S$ 各纤维上的 $r$-循环族之间诱导双射。这些限制映射
$Z_r(X_s) \to Z_r(U_s)$ 与基变换和特化相容；见引理
\ref{relative-cycles-lemma-compatibilities} 和 \ref{relative-cycles-lemma-specialization-flat-pullback}。
引理由此容易得出；细节略。
\end{proof}

\begin{lemma}
\label{relative-cycles-lemma-seminormalize-base}
设 $g : S' \to S$ 为局部 Noether 概形之间诱导剩余域同构的泛同胚。设
$f : X \to S$ 局部有限型，并令 $X' = S' \times_S X$。设 $r \geq 0$。
则沿 $g$ 的基变换确定双射 $z(X/S, r) \to z(X'/S', r)$。
\end{lemma}

\begin{proof}
由引理 \ref{relative-cycles-lemma-pullback-universally-bijective}，$X/S$ 各纤维上的
$r$-循环族之群与 $X'/S'$ 各纤维上的 $r$-循环族之群之间存在双射。设
$\alpha$ 为 $X/S$ 各纤维上的一个 $r$-循环族，$\alpha' = g^*\alpha$ 为其
基变换。若 $R$ 是离散赋值环，则任意态射 $h : \Spec(R) \to S$ 都唯一地
分解为 $g \circ h'$，其中 $h' : \Spec(R) \to S'$。事实上，态射
$S' \times_S \Spec(R) \to \Spec(R)$ 是在剩余域上诱导双射的泛同胚，
因而有截面（例如因为 $R$ 是半正规环；见“态射”第
\ref{morphisms-section-seminormalization} 节）。因此，$\alpha$ 与特化相容
（即为相对 $r$-循环）这一条件，等价于 $\alpha'$ 与特化相容。
\end{proof}






\section{等维相对循环}
\label{relative-cycles-section-equidimensional}

\noindent
定义如下。

\begin{definition}
\label{relative-cycles-definition-equidimensional}
设 $f : X \to S$ 为概形态射。假设 $S$ 局部 Noether，且 $f$ 局部有限型。
设 $r \geq 0$ 为整数。如果 $X/S$ 上相对 $r$-循环 $\alpha$ 的支撑
（评注 \ref{relative-cycles-remark-supports-family}）包含在闭子集 $W \subset X$ 中，且后者
在 $S$ 上的相对维数 $\leq r$，则称 $\alpha$ {\it 等维}。
$X/S$ 上所有等维相对 $r$-循环所成的群记作 $z_{equi}(X/S, r)$。
\end{definition}

\begin{example}
\label{relative-cycles-example-not-equidimensional}
\begin{reference}
\cite[Example 3.1.9]{SV}
\end{reference}
存在并非等维的相对 $r$-循环。具体地，设 $k$ 为域，并令
$X = \Spec(k[x, y, t])$ 位于 $S = \Spec(k[x, y])$ 之上。设 $s$ 为 $S$
的一点，并以 $a, b \in \kappa(s)$ 表示 $x$ 和 $y$ 的像。考虑如下定义的
$X/S$ 各纤维上的 $0$-循环族 $\alpha$：
\begin{enumerate}
\item 若 $b = 0$，则 $\alpha_s = 0$；否则
\item $\alpha_s = [p] - [q]$，其中 $p$（相应地，$q$）是
$\Spec(\kappa(s)[t])$ 的 $\kappa(s)$-有理点，满足 $t = a/b$
（相应地，$t = (a + b^2)/b$）。
\end{enumerate}
留给读者证明它与特化相容；思路是，对 $\kappa(s)$ 上任意满足
$v(b) > 0$ 的赋值 $v$，$a/b$ 与 $(a + b^2)/b = a/b + b$ 在其剩余域上的
$\mathbf{P}^1$ 中趋于同一点。另一方面，$\alpha$ 的支撑之闭包包含
$(0, 0)$ 上的整条纤维。
\end{example}

\begin{lemma}
\label{relative-cycles-lemma-equidimensional-functoriality}
设 $f : X \to S$ 为概形态射。假设 $S$ 局部 Noether，且 $f$ 局部有限型。
设 $r \geq 0$ 为整数，并令 $\alpha$ 为 $X/S$ 上的相对 $r$-循环。若
$\alpha$ 等维，则 $\alpha$ 的任意限制、基变换或平坦拉回都等维。
\end{lemma}

\begin{proof}
略。
\end{proof}

\begin{lemma}
\label{relative-cycles-lemma-check-equidimensional}
设 $f : X \to S$ 为概形态射。假设 $S$ 局部 Noether，且 $f$ 局部有限型。
设 $r \geq 0$ 为整数，并令 $\alpha$ 为 $X/S$ 上的相对 $r$-循环。则检验
$\alpha$ 等维时，可以在 $X$ 和 $S$ 上 Zariski 局部地进行。
\end{lemma}

\begin{proof}
确实，$\alpha$ 等维这一条件恰好意味着 $\alpha$ 的支撑之闭包在 $S$ 上的
相对维数 $\leq r$。由于取闭包与限制到开集可交换，引理得证
（略去一个小细节）。
\end{proof}

\begin{lemma}
\label{relative-cycles-lemma-equidimensional-fppf-descent}
设 $f : X \to S$ 为概形态射。假设 $S$ 局部 Noether，且 $f$ 局部有限型。
设 $r \geq 0$ 为整数，并令 $\alpha$ 为 $X/S$ 上的相对 $r$-循环。设
$\{g_i : S_i \to S\}$ 为 fppf-覆盖。则 $\alpha$ 等维，当且仅当每个
基变换 $g_i^*\alpha$ 都等维。
\end{lemma}

\begin{proof}
若 $\alpha$ 等维，则由引理 \ref{relative-cycles-lemma-equidimensional-functoriality}，
每个 $g_i^*\alpha$ 也等维。反之，假设每个 $g_i^*\alpha$ 都等维。以 $W$
表示 $\text{Supp}(\alpha)$ 在 $X$ 中的闭包。由于
$g_i : S_i \to S$ 泛开（因为它平坦且局部有限表示），态射
$f_i : X_i = S_i \times_S X \to X$ 也泛开。记 $\alpha_i = g_i^*\alpha$。
由引理 \ref{relative-cycles-lemma-support-family}，
$\text{Supp}(\alpha_i) = f_i^{-1}(\text{Supp}(\alpha))$。由于 $f_i$ 开，
$W_i = f_i^{-1}(W)$ 是 $\text{Supp}(\alpha_i)$ 的闭包。因此根据假设，
态射 $W_i \to S_i$ 的相对维数 $\leq r$。由“态射”引理
\ref{morphisms-lemma-dimension-fibre-after-base-change}（以及态射
$S_i \to S$ 联合满射这一事实），可知 $W \to S$ 的相对维数 $\leq r$。
\end{proof}

\begin{lemma}
\label{relative-cycles-lemma-equidimensional-descent-pullbacks}
设 $f : X \to S$ 为概形态射。假设 $S$ 局部 Noether，且 $f$ 局部有限型。
设 $r, e \geq 0$ 为整数，并令 $\alpha$ 为 $X/S$ 上的相对 $r$-循环。
设 $\{f_i : X_i \to X\}$ 为一族联合满射的局部有限型平坦态射，且其相对
维数为 $e$。则 $\alpha$ 等维，当且仅当每个平坦拉回 $f_i^*\alpha$ 都等维。
\end{lemma}

\begin{proof}
略。提示：仿照引理 \ref{relative-cycles-lemma-equidimensional-fppf-descent} 的证明，
可证 $\alpha$ 的支撑之闭包 $W$ 沿 $f_i$ 的逆像，是 $f_i^*\alpha$ 的支撑之
闭包 $W_i$。于是，若对所有 $i$，$W_i \to S$ 的相对维数
$\leq r + e$，则 $W \to S$ 的相对维数 $\leq r$。
\end{proof}





\section{权函数与相对零循环}
\label{relative-cycles-section-weightings}

\noindent
本节说明，对拟有限态射，相对零循环与“态射续篇”第
\ref{more-morphisms-section-weightings} 节所定义的权函数密切相关。

\medskip\noindent
设 $S$ 为局部 Noether 概形，$f : X \to S$ 为概形的局部拟有限态射。则
$z(X/S, 0) = z_{equi}(X/S, 0)$，并且当 $r > 0$ 时 $z(X/S, r) = 0$。
给定 $\alpha \in z(X/S, 0)$，定义映射
$$
w_\alpha : X \longrightarrow \mathbf{Z},\quad
x \mapsto \alpha(x) [\kappa(x) : \kappa(s)]_i \quad\text{其中 }s = f(x)
$$
这里 $\alpha(x)$ 表示纤维 $X_s$ 上的 $0$-循环 $\alpha_s$ 中 $x$ 的系数，
而 $[K : k]_i$ 表示有限域扩张的不可分次数。

\begin{lemma}
\label{relative-cycles-lemma-weightings-pre}
在上述情形中，若 $g : S' \to S$ 为局部 Noether 概形之间的态射，则
$w_{g^*\alpha} = w_\alpha \circ g'$，其中 $g' : X' \to X$ 是投影
$X' = S' \times_S X \to X$。
\end{lemma}

\begin{proof}
设 $x' \in X'$ 在 $S', S, X$ 中的像分别为 $s', s, x$。则沿
$\kappa(s')/\kappa(s)$ 对 $[x]$ 作基变换后，$[x']$ 的系数等于局部环
$(\kappa(s') \otimes_{\kappa(s)} \kappa(x))_\mathfrak q$ 的长度；这里
$\mathfrak q$ 是与 $x'$ 对应的素理想。因此，只需证明
$$
[\kappa(x) : \kappa(s)]_i =
\text{length}((\kappa(s') \otimes_{\kappa(s)} \kappa(x))_\mathfrak q)
[\kappa(x') : \kappa(s')]_i
$$
即可得到与基变换的相容性。设 $k/\kappa(s')$ 为一个代数闭包。选取位于
$\mathfrak q$ 之上的素理想
$\mathfrak p \subset k \otimes_{\kappa(s)} \kappa(x)$。假设能够证明
$$
[\kappa(x) : \kappa(s)]_i =
\text{length}((k \otimes_{\kappa(s)} \kappa(x))_\mathfrak p)
\quad\text{且}\quad
[\kappa(x') : \kappa(s')]_i =
\text{length}((k \otimes_{\kappa(s')} \kappa(x'))_\mathfrak p)
$$
则结论成立，因为
$$
\text{length}((\kappa(s') \otimes_{\kappa(s)} \kappa(x))_\mathfrak q)
\text{length}((k \otimes_{\kappa(s')} \kappa(x'))_\mathfrak p)
=
\text{length}((k \otimes_{\kappa(s)} \kappa(x))_\mathfrak p)
$$
由“代数”引理 \ref{algebra-lemma-pullback-module} 以及
$\kappa(s') \otimes_{\kappa(s)} \kappa(x) \to k \otimes_{\kappa(s)} \kappa(x)$
的平坦性得出。为证明上述两个等式，只需证明第一个。令
$\kappa(x)/\kappa/\kappa(s)$ 为“域”引理
\ref{fields-lemma-separable-first} 中构造的子域。于是
$$
k \otimes_{\kappa(s)} \kappa(x) =
\prod\nolimits_{\sigma : \kappa \to k}
k \otimes_{\sigma, \kappa} \kappa(x)
$$
且每个因子都是次数为
$[\kappa(x) : \kappa] = [\kappa(x) : \kappa(s)]_i$ 的局部环，正合所需。
\end{proof}

\begin{lemma}
\label{relative-cycles-lemma-weightings}
设 $S$ 为局部 Noether 概形，$f : X \to S$ 为概形的局部拟有限态射。
\begin{enumerate}
\item 对 $\alpha \in z(X/S, 0)$，上述构造的映射
$w_\alpha : X \to \mathbf{Z}$ 是权函数。
\item 若 $X$ 拟紧，则给定权函数 $w : X \to \mathbf{Z}$，存在整数
$n > 0$，使得对某个 $\alpha \in z(X/S, 0)$ 有 $nw = w_\alpha$。
\item 若 $S$ 是 $\mathbf{F}_p$ 上的概形，则 (2) 中的整数 $n$ 可以取为
素数 $p$ 的幂。
\end{enumerate}
\end{lemma}

\begin{proof}
设 $\alpha \in z(X/S, 0)$，并选取图
$$
\xymatrix{
X \ar[d]_f & U \ar[l]^h \ar[d]^\pi \\
Y & V \ar[l]_g
}
$$
如“态射续篇”定义 \ref{more-morphisms-definition-weighting} 所述。以
$\beta \in z(U/V, 0)$ 表示基变换 $g^*\alpha$ 的限制。由与基变换的相容性
（引理 \ref{relative-cycles-lemma-weightings-pre}），$w_\beta = w_\alpha \circ h$，所以只需
证明 $\int_\pi w_\beta$ 在 $V$ 上局部常值。注意
\begin{align*}
\left( \int_\pi w_\beta \right)(v)
& =
\sum\nolimits_{u \in U, \pi(u) = v} 
\beta(u) [\kappa(u) : \kappa(v)]_i [\kappa(u) : \kappa(v)]_s \\
& =
\sum\nolimits_{u \in U, \pi(u) = v} \beta(u)[\kappa(u) : \kappa(v)]
\end{align*}
最后一个表达式是 $\pi_*\beta \in z(V/V, 0)$ 中 $v$ 的系数。由引理
\ref{relative-cycles-lemma-relative-cycle-smooth}，该函数在 $V$ 上局部常值。

\medskip\noindent
反之，设 $w : X \to S$ 为权函数，且 $X$ 拟紧。选取充分可除的整数 $n$。
令 $\alpha$ 为 $X/S$ 各纤维上的 $0$-循环族，使得对 $s \in S$ 有
$$
\alpha_s =
\sum\nolimits_{f(x) = s} \frac{n w(x)}{[\kappa(x) : \kappa(s)]_i} [x]
$$
作为 $X_s$ 上的零循环。这一定义有意义，因为 $f$ 的纤维普遍有界
（“态射”引理
\ref{morphisms-lemma-locally-quasi-finite-qc-source-universally-bounded}），
故可以选取 $n$，使右边对所有 $s \in S$ 都取整数。只要证明 $\alpha$
是相对 $0$-循环，引理的最后一个断言也随之成立。为此，需要证明
$\alpha$ 与沿离散赋值环的特化相容。注意，$w_\alpha$ 的构造以及引理
\ref{relative-cycles-lemma-weightings-pre} 的证明都没有使用 $\alpha$ 是相对循环这一点。
因此 $w_\alpha = nw$，且该等式在基变换后仍成立。此外，权函数的基变换
仍是权函数；见“态射续篇”引理
\ref{more-morphisms-lemma-weighting-base-change}。从而归约到下一段所述问题。

\medskip\noindent
假设 $S$ 是离散赋值环的谱，其一般点为 $\eta$、闭点为 $0$。设
$w : X \to S$ 为权函数，且 $X$ 在 $S$ 上拟有限。令 $\alpha$ 为上一段
构造的 $X/S$ 各纤维上的 $0$-循环族（取适当的 $n$）。需要证明
$sp_{X/S}(\alpha_\eta) = \alpha_0$。令 $\beta \in z(X/S, 0)$ 为 $X/S$
上的相对 $0$-循环，满足 $\beta_\eta = \alpha_\eta$ 且
$\beta_0 = sp_{X/S}(\alpha_\eta)$。于是
$w' = w_\beta - nw : X \to \mathbf{Z}$ 是权函数（使用上面的结论），
并在 $X$ 中映到 $\eta$ 的点处为零。由“态射续篇”引理
\ref{more-morphisms-lemma-weighting-specialization}，$w' = 0$。这意味着
$w_\beta = nw$，故如所需，$\alpha = \beta$。
\end{proof}

\begin{example}
\label{relative-cycles-example-n-must-be-positive-sometimes}
设 $p$ 为素数，$k'/k$ 为次数 $p^f$ 的纯不可分域扩张。令
$X = \Spec(k')$、$S = \Spec(k)$。把 $X$ 的唯一一点 $x$ 映到 $1$ 的映射
$w : X \to \mathbf{Z}$ 是 $X \to S$ 的权函数。另一方面，群
$z(X/S, 0)$ 由循环 $\alpha = [x]$ 自由生成。因此在此情形中，引理
\ref{relative-cycles-lemma-weightings} 所需的最小整数 $n$ 是 $n = p^f$。
\end{example}

\begin{example}
\label{relative-cycles-example-Merkurjev}
上述讨论可以用来“解释” A.S.~Merkurjev 的一个例子
\cite[Example 3.5.10]{SV}。设 $p$ 为素数，并令
$k = \mathbf{F}_p(a, b)$ 为 $\mathbf{F}_p$ 关于 $a$ 和 $b$ 的纯超越扩张。
令
$$
A = k[x, y, z]/(a x^p + b y^p - z^p)
$$
这是 $k$ 上有限型的正规整环。谱 $S = \Spec(A)$ 除了唯一的奇点 $s$ 外
正则；该点对应于 $A$ 的极大理想 $(x, y, z)$，其剩余域为 $k$。令
$k' = k(a^{1/p}, b^{1/p})$、$B = k'[x, y]$。利用嵌入 $k \to k'$，考虑
映射 $A \to B$：它把 $x$ 映到 $x$、把 $y$ 映到 $y$、把 $z$ 映到
$a^{1/p}x + b^{1/p}y$。令 $X = \Spec(B)$，并考虑与环映射 $A \to B$
对应的态射 $f : X \to S$。由“态射续篇”引理
\ref{more-morphisms-lemma-weighting-quasi-finite-Noetherian} 为 $f$ 产生的
权函数 $w : X \to \mathbf{Z}$ 恒等于 $1$，因为 $B$ 的分式域在 $A$ 上
纯不可分！令 $V = f^{-1}(U) \subset X$。由奇迹平坦性，$V \to U$ 平坦，
次数为 $p$。因此循环 $\alpha = [V/V/U]_0 \in z(V/U, 0)$ 满足
$w_\alpha = pw|_V$。然而，位于 $s$ 上方的 $X$ 的唯一一点 $x$ 的剩余域
为 $k'$，它在 $k$ 上的次数为 $p^2$；故不存在 $z(X/S, 0)$ 中限制为
$\alpha$ 的元素 $\beta$。事实上，由“态射续篇”引理
\ref{more-morphisms-lemma-weighting-specialization}，权函数
$pw - w_\beta$ 必为零；这将由上面引理证明中的公式蕴含
$$
\beta_s = \frac{pw(x)}{[\kappa(x) : \kappa(s)]_i} [x] =
\frac{1}{p}[x]
$$
于 $Z_0(X_s)$ 中，而这是不可能的\footnote{在 \cite{SV} 中，确实存在
相对循环群 $Cycl(X/S, 0)$ 中与 $pw$ 对应的元素，而且该循环确实特化为
$X_s$ 上系数非整数的循环。}。当然，循环 $p\alpha$ 是 $z(X/S, 0)$ 中
唯一一个元素的限制。
\end{example}









\section{有效相对循环}
\label{relative-cycles-section-effective}


\noindent
定义如下。

\begin{definition}
\label{relative-cycles-definition-effective}
设 $f : X \to S$ 为概形态射。假设 $S$ 局部 Noether，且 $f$ 局部有限型。
设 $r \geq 0$ 为整数。若对所有 $s \in S$，$\alpha_s$ 都是有效循环
（Chow 同调，定义 \ref{chow-definition-effective-cycle}），则称 $X/S$ 上的
相对 $r$-循环 $\alpha$ {\it 有效}。$X/S$ 上所有有效相对 $r$-循环所成的
幺半群记作 $z^{eff}(X/S, r)$。
\end{definition}

\noindent
下面将证明有效相对循环等维；见引理
\ref{relative-cycles-lemma-effective-equidimensional}。

\begin{lemma}
\label{relative-cycles-lemma-effective-functoriality}
设 $f : X \to S$ 为概形态射。假设 $S$ 局部 Noether，且 $f$ 局部有限型。
设 $r \geq 0$ 为整数，并令 $\alpha$ 为 $X/S$ 上的相对 $r$-循环。若
$\alpha$ 有效，则 $\alpha$ 的任意限制、基变换、平坦拉回或固有前推都有效。
\end{lemma}

\begin{proof}
略。
\end{proof}

\begin{lemma}
\label{relative-cycles-lemma-check-effective}
设 $f : X \to S$ 为概形态射。假设 $S$ 局部 Noether，且 $f$ 局部有限型。
设 $r \geq 0$ 为整数，并令 $\alpha$ 为 $X/S$ 上的相对 $r$-循环。则检验
$\alpha$ 有效时，可以在 $X$ 和 $S$ 上 Zariski 局部地进行。
\end{lemma}

\begin{proof}
略。
\end{proof}

\begin{lemma}
\label{relative-cycles-lemma-effective-descent}
设 $f : X \to S$ 为概形态射。假设 $S$ 局部 Noether，且 $f$ 局部有限型。
设 $r \geq 0$ 为整数，并令 $\alpha$ 为 $X/S$ 上的相对 $r$-循环。设
$g : S' \to S$ 为满射。则 $\alpha$ 有效，当且仅当基变换 $g^*\alpha$
有效。
\end{lemma}

\begin{proof}
略。
\end{proof}

\begin{lemma}
\label{relative-cycles-lemma-effective-descent-pullbacks}
设 $f : X \to S$ 为概形态射。假设 $S$ 局部 Noether，且 $f$ 局部有限型。
设 $r, e \geq 0$ 为整数，并令 $\alpha$ 为 $X/S$ 上的相对 $r$-循环。
设 $\{f_i : X_i \to X\}$ 为一族联合满射的局部有限型平坦态射，且其相对
维数为 $e$。则 $\alpha$ 有效，当且仅当每个平坦拉回 $f_i^*\alpha$ 都有效。
\end{lemma}

\begin{proof}
略。
\end{proof}

\begin{lemma}
\label{relative-cycles-lemma-effective-support-closed}
设 $f : X \to S$ 为概形态射。假设 $S$ 局部 Noether，且 $f$ 局部有限型。
设 $r, e \geq 0$ 为整数，并令 $\alpha$ 为 $X/S$ 上的相对 $r$-循环。
若 $\alpha$ 有效，则 $\text{Supp}(\alpha)$ 在 $X$ 中闭。
\end{lemma}

\begin{proof}
设 $g : S' \to S$ 是一个不可约分支的浸入，并把该分支视作整闭子概形。
由引理 \ref{relative-cycles-lemma-effective-functoriality} 和
\ref{relative-cycles-lemma-support-family}，只需证明基变换 $g^*\alpha$ 的支撑在
$S' \times_S S$ 中闭。因此可以假设 $S$ 为一般点是 $\eta$ 的整概形。
我们将证明 $\text{Supp}(\alpha)$ 是 $\text{Supp}(\alpha_\eta)$ 的闭包。
为此，任取 $s \in S$。存在态射 $g : S' \to S$，其中 $S'$ 是离散赋值环
的谱，它把一般点 $\eta' \in S'$ 映到 $\eta$，把闭点 $0 \in S'$ 映到
$s$；见“性质”引理
\ref{properties-lemma-locally-Noetherian-specialization-dvr}。于是只需证明
$g^*\alpha$ 的支撑等于 $\text{Supp}((g^\alpha)_{\eta'})$ 的闭包。这就
归约到下一段所述情形。

\medskip\noindent
此处 $S$ 是离散赋值环的谱，其一般点为 $\eta$、闭点为 $0$。需要证明
$\text{Supp}(\alpha)$ 是 $\text{Supp}(\alpha_\eta)$ 的闭包。由于 $\alpha$
有效，可以写成 $\alpha_\eta = \sum n_i[Z_i]$，其中 $n_i > 0$，且
$Z_i \subset X_\eta$ 是维数为 $r$ 的整闭子概形。由于
$\alpha_0 = sp_{X/S}(\alpha_\eta)$，有
$\alpha_0 = \sum n_i [\overline{Z}_{i, 0}]_r$，其中 $\overline{Z}_i$ 是
$Z_i$ 的闭包。由“代数簇”引理
\ref{varieties-lemma-dominate-valuation-ring-dimension-fibres}，
$\overline{Z}_{i, 0}$ 是维数为 $r$ 的等维概形。由于 $n_i > 0$，
$\text{Supp}(\alpha_0)$ 等于所有 $\overline{Z}_{i, 0}$ 的并；它就是
$\bigcup \overline{Z}_i$ 在 $0$ 上的纤维，而后者又是 $\bigcup Z_i$ 的
闭包，正合所需。
\end{proof}

\begin{lemma}
\label{relative-cycles-lemma-effective-equidimensional}
设 $f : X \to S$ 为概形态射。假设 $S$ 局部 Noether，且 $f$ 局部有限型。
设 $r, e \geq 0$ 为整数，并令 $\alpha$ 为 $X/S$ 上的相对 $r$-循环。
若 $\alpha$ 有效，则 $\alpha$ 等维。
\end{lemma}

\begin{proof}
假设 $\alpha$ 有效。由引理 \ref{relative-cycles-lemma-effective-support-closed}，支撑
$\text{Supp}(\alpha)$ 在 $X$ 中闭。由于 $\text{Supp}(\alpha) \to S$ 的
纤维是循环 $\alpha_s$ 的支撑，因而维数为 $r$，所以 $\alpha$ 等维。
\end{proof}

\begin{remark}
\label{relative-cycles-remark-representable}
设 $f : X \to S$ 为概形态射，其中 $S$ 局部 Noether，且 $f$ 局部有限型。
我们可以问，下述反变函子
$$
\begin{matrix}
\text{局部有限型的概形 }S'\text{，} \\
\text{其底概形为 }S
\end{matrix}
\longrightarrow
z^{eff}(X'/S', r)\text{，其中 }X' = S' \times_S X
$$
是否可表。由于 $z(X'/S', r) = z(X'_{red}/S'_{red}, r)$，答案不可能为是
（我们把构造一个实际反例留给读者）。更合适的问题是：能否在左端找到一个
子范畴，使该函子限制在其上后可表。引理 \ref{relative-cycles-lemma-seminormalize-base}
提示我们至少应限制到 $S$ 上半正规概形的范畴。

\medskip\noindent
若 $S/\Spec(\mathbf{Q})$ 是 Nagata 的，且 $f$ 为投射态射，则事实证明，
$S' \mapsto z^{eff}(X'/S', r)$ 在半正规的 $S'$ 所成范畴上可表。
粗略地说，这正是 \cite[Theorem 3.21]{KRC} 的内容。

\medskip\noindent
若 $S$ 有正特征点，那么即使把半正规性换成弱正规性，上述结论也不再成立；
若对任意双有理泛同胚 $T' \to T$ 都存在截面，则称局部 Noether 概形 $T$
为弱正规的。一个例子是取 $S = \Spec(k)$，其中 $k$ 为特征 $2$ 的域，并考虑
$X = \mathbf{A}^2_k$ 上次数为 $2$ 的 $0$-循环。具体而言，在
$W = X \times_S X$ 上有一个典范相对 $0$-循环
$\alpha \in z^{eff}(X_W/W, 0)$：对 $w = (x_1, x_2) \in W = X^2$，有
$\alpha_w = [x_1] + [x_2]$。此循环在交换两因子的对合
$\sigma : W \to W$ 下不变。由于 $W$ 光滑（从而正规，进而弱正规），若
$z(-/-, r)$ 在 $k$ 上有限型弱正规概形的范畴上由 $M$ 表示，我们便得到从
$W$ 到 $M$ 的一个 $\sigma$-不变态射。它继而定义从商概形
$\text{Sym}^2_S(X) = W/\langle \sigma \rangle$ 到 $M$ 的态射。由于
$\text{Sym}^2_S(X)$ 正规，由 $M$ 的模性质，我们会在
$X \times_S \text{Sym}^2_S(X) / \text{Sym}^2_S(X)$ 上得到一个相对
$0$-循环 $\beta$，其到 $W$ 的拉回为 $\alpha$。然而，这样的循环 $\beta$
并不存在。事实上，写 $X = \Spec(k[u, v])$，则概形
$\text{Sym}^2_S(X)$ 是下述环的谱：
$$
k[u_1 + u_2, u_1u_2, v_1 + v_2, v_1v_2, u_1v_1 + u_2v_2]
\subset
k[u_1, u_2, v_1, v_2]
$$
对角线 $u_1 = u_2, v_1 = v_2$ 在 $\text{Sym}^2_S(X)$ 中的像是闭子概形
$V = \Spec(k[u_1^2, v_1^2])$；这里用到了 $k$ 的特征为 $2$。考察 $V$ 的
一般点 $\eta$，则 $\beta_\eta$ 会是
$\mathbf{A}^2_{k(u_1^2, v_1^2)}$ 上次数为 $2$ 的零维循环，而它到
$\mathbf{A}^2_{k(u_1, u_2)}$ 的拉回将是
$2[\text{坐标为 } (u_1, v_2)]$。这显然不可能。

\medskip\noindent
上述讨论并不与 \cite[Theorem 4.13]{KRC} 矛盾，因为该定理中的 Chow 簇只对
一个函子给出粗表示（事实上是两个不同的函子；经过一些工作可以看出，在
$X$ 投射时只有其中一个与我们的函子一致）。类似地，
\cite[Section 4.4]{SV} 证明：当 $X/S$ 投射时，预层
$S' \mapsto z^{eff}(S' \times_S X/S', r)$ 的 $h$-层化等于某个可表函子的
$h$-层化。
\end{remark}

\begin{remark}
\label{relative-cycles-remark-equidimensional-over-geometrically-unibranch}
设 $f : X \to S$ 为概形态射，$r \geq 0$，并设 $Z \subset X$ 为闭子概形。
假设
\begin{enumerate}
\item $S$ 是 Noether 且几何单枝的，
\item $f$ 有限型，并且
\item $Z \to S$ 的相对维数 $\leq r$。
\end{enumerate}
则对所有充分可除的整数 $n \geq 1$，存在唯一的 $X/S$ 上有效相对
$r$-循环 $\alpha$，使得对 $S$ 的每个一般点 $\eta$ 都有
$\alpha_\eta = n[Z_\eta]_r$。这是 \cite[Theorem 3.4.2]{SV} 的一种改述。
若以后需要这一结果，我们将在此准确陈述并证明它。
\end{remark}











\section{固有相对循环}
\label{relative-cycles-section-proper}

\noindent
在我们的设定中，下面大概是恰当的定义。

\begin{definition}
\label{relative-cycles-definition-proper}
设 $f : X \to S$ 为概形态射。假设 $S$ 局部 Noether，且 $f$ 局部有限型。
设 $r \geq 0$ 为整数。若 $\alpha$ 的支撑（注记
\ref{relative-cycles-remark-supports-family}）包含在某个在 $S$ 上固有的闭子集
$W \subset X$ 中（概形上同调，定义
\ref{coherent-definition-proper-over-base}），则称 $X/S$ 上的相对
$r$-循环 $\alpha$ 为{\it 固有相对循环}。$X/S$ 上所有固有相对
$r$-循环所成的群记为 $c(X/S, r)$。
\end{definition}

\noindent
由概形上同调，引理 \ref{coherent-lemma-closed-closed-proper-over-base}，
这恰好意味着支撑的闭包在底上固有。要看出这些循环构成一个群，可用
概形上同调，引理 \ref{coherent-lemma-union-closed-proper-over-base}。

\begin{lemma}
\label{relative-cycles-lemma-proper-functoriality}
设 $f : X \to S$ 为概形态射。假设 $S$ 局部 Noether，且 $f$ 局部有限型。
设 $r \geq 0$ 为整数，并令 $\alpha$ 为 $X/S$ 上的相对 $r$-循环。若
$\alpha$ 固有，则 $\alpha$ 的任意基变换都固有。
\end{lemma}

\begin{proof}
略。
\end{proof}

\begin{lemma}
\label{relative-cycles-lemma-proper-h-descent}
设 $f : X \to S$ 为概形态射。假设 $S$ 局部 Noether，且 $f$ 局部有限型。
设 $r \geq 0$ 为整数，令 $\alpha$ 为 $X/S$ 上的相对 $r$-循环，并设
$\{g_i : S_i \to S\}$ 为一个 h-覆盖。则 $\alpha$ 固有，当且仅当每个
基变换 $g_i^*\alpha$ 都固有。
\end{lemma}

\begin{proof}
若 $\alpha$ 固有，则由引理 \ref{relative-cycles-lemma-proper-functoriality}，每个
$g_i^*\alpha$ 也固有。反之，假设每个 $g_i^*\alpha$ 都固有。为证明
$\alpha$ 固有，显然只需在 $S$ 上仿射局部地工作，故可设 $S$ 仿射。
于是可用一族 $\{T_j \to S\}$ 加细覆盖 $\{S_i \to S\}$，其中
$g : T \to S$ 是固有满态射，且 $T = \bigcup T_j$ 为开覆盖。
因此，$\beta = g^*\alpha$ 在 $Y = T \times_S X$ 上作为 $T$ 上的循环是
固有的。由引理 \ref{relative-cycles-lemma-support-family}，$\beta$ 的支撑是 $\alpha$ 的
支撑在态射 $f : Y \to X$ 下的逆像。因此，$f^{-1}\text{Supp}(\alpha)$
在该概形中的闭包 $W \subset Y$ 在 $T$ 上固有。又因态射 $T \to S$
固有，所以 $W$ 在 $S$ 上固有。于是由概形上同调，引理
\ref{coherent-lemma-functoriality-closed-proper-over-base}，像
$f(W) \subset X$ 是在 $S$ 上固有的闭子集。由于 $f(W)$ 包含
$\text{Supp}(\alpha)$，故 $\alpha$ 固有。
\end{proof}



\section{固有且等维的相对循环}
\label{relative-cycles-section-proper-equidimensional}

\noindent
设 $f : X \to S$ 为概形态射。假设 $S$ 局部 Noether，且 $f$ 局部有限型。
设 $r \geq 0$ 为整数。若 $X/S$ 上的相对 $r$-循环 $\alpha$ 既等维
（定义 \ref{relative-cycles-definition-equidimensional}）又固有（定义
\ref{relative-cycles-definition-proper}），则称 $\alpha$ 为{\it 固有且等维的相对循环}。
$X/S$ 上所有固有且等维的相对 $r$-循环所成的群记为
$c_{equi}(X/S, r)$。

\medskip\noindent
类似地，若 $X/S$ 上的相对 $r$-循环 $\alpha$ 既有效（定义
\ref{relative-cycles-definition-effective}）又固有（定义 \ref{relative-cycles-definition-proper}），
则称 $\alpha$ 为{\it 固有且有效的相对循环}。$X/S$ 上所有固有且有效的相对
$r$-循环所成的幺半群记为 $c^{eff}(X/S, r)$。注意，由引理
\ref{relative-cycles-lemma-effective-equidimensional}，这些循环都是等维的。

\medskip\noindent
于是有如下包含映射图：
$$
\xymatrix{
c^{eff}(X/S, r) \ar[r] \ar[d] &
c_{equi}(X/S, r) \ar[r] \ar[d] &
c(X/S, r) \ar[d] \\
z^{eff}(X/S, r) \ar[r] &
z_{equi}(X/S, r) \ar[r] &
z(X/S, r)
}
$$



\section{循环上的作用}
\label{relative-cycles-section-action}

\noindent
设 $S$ 为局部 Noether、普遍链状的概形，并赋予维数函数 $\delta$，参见
Chow 同调，第 \ref{chow-section-setup} 节。设 $X \to Y$ 为 $S$ 上的
概形态射，且二者均在 $S$ 上局部有限型。设 $r \geq 0$。最后，令
$\alpha$ 为 $X/Y$ 的纤维上的 $r$-循环族。对 $e \in \mathbf{Z}$，
我们将构造运算
$$
\alpha \cap - : Z_e(Y) \longrightarrow Z_{r + e}(X)
$$
具体而言，给定 $\beta \in Z_e(Y)$，写成 $\beta = \sum n_i[Z_i]$，其中
$Z_i \subset Y$ 是 $\delta$-维数为 $e$ 的整闭子概形，且族 $Z_i$ 在概形
$Y$ 中局部有限。令 $y_i \in Z_i$ 为一般点，并写
$\alpha_{y_i} = \sum m_{ij} [V_{ij}]$。于是 $V_{ij} \subset X_{y_i}$
是维数为 $r$ 的整闭子概形，且族 $V_{ij}$ 在概形 $X_{y_i}$ 中局部有限。
定义
$$
\alpha \cap \beta = \sum n_i m_{ij} [\overline{V}_{ij}]
\quad\in\quad
Z_{r + e}(X)
$$
这里，$\overline{V}_{ij} \subset X$ 是态射 $V_{ij} \to X_{y_i} \to X$
的概形论像；等价地，$\overline{V}_{ij} \subset X$ 是一个支配地映到
$Z_i \subset Y$、且一般纤维为 $V_{ij}$ 的整闭子概形。容易推出
$\dim_\delta(\overline{V}_{ij}) = r + e$，并且闭子概形族
$\overline{V}_{ij} \subset X$ 局部有限（验证从略）。因此
$\alpha \cap \beta$ 的确是 $Z_{r + e}(X)$ 的元素。

\begin{lemma}
\label{relative-cycles-lemma-action-bilinear}
上述构造是双线性的，即
$(\alpha_1 + \alpha_2) \cap \beta = \alpha_1 \cap \beta +
\alpha_2 \cap \beta$，且 $\alpha \cap (\beta_1 + \beta_2) =
\alpha \cap \beta_1 + \alpha \cap \beta_2$.
\end{lemma}

\begin{proof}
略。
\end{proof}

\begin{lemma}
\label{relative-cycles-lemma-action-opens}
若 $U \subset X$ 与 $V \subset Y$ 为开集，且 $f(U) \subset V$，则
$(\alpha \cap \beta)|_U$ 等于 $\alpha|_U \cap \beta|_V$。
\end{lemma}

\begin{proof}
由上面对 $\alpha \cap \beta$ 的显式描述立即可得。
\end{proof}

\begin{lemma}
\label{relative-cycles-lemma-action-base-change}
构造 $\alpha \cap \beta$ 与平坦基变换和平坦拉回相容（详见证明）。
\end{lemma}

\begin{proof}
令 $(S, \delta)$、$(S', \delta')$、$g : S' \to S$ 与
$c \in \mathbf{Z}$ 如 Chow 同调，情形 \ref{chow-situation-setup-base-change}
所述。设 $X \to Y$ 为 $S$ 上局部有限型的概形态射，记 $X' \to Y'$
为 $X \to Y$ 沿 $g$ 的基变换。令 $\alpha$ 为 $X/Y$ 的纤维上的
$r$-循环族，并令 $\beta \in Z_e(Y)$。记 $\alpha'$ 为 $\alpha$ 沿
$Y' \to Y$ 的基变换，记 $\beta' = g^*\beta \in Z_{e + c}(Y')$ 为
$\beta$ 沿 $g$ 的拉回，参见 Chow 同调，第 \ref{chow-section-change-base}
节。与基变换相容，是指 $\alpha' \cap \beta'$ 是
$\alpha \cap \beta$ 的基变换。

\medskip\noindent
先证与基变换相容。由于要证明的是 $X'$ 上循环的等式，由引理
\ref{relative-cycles-lemma-action-opens}，可以在 $Y$ 上局部地工作。因此可设 $Y$ 仿射。
特别地，$\beta$ 是素循环的有限线性组合。由于 $- \cap -$ 关于第二个
变量线性（引理 \ref{relative-cycles-lemma-action-bilinear}），只需在
$\beta = [Z]$ 时证明该等式，其中 $Z \subset Y$ 是 $\delta$-维数为
$e$ 的整闭子概形。

\medskip\noindent
令 $y \in Z$ 为一般点。写 $\alpha_y = \sum m_j [V_j]$，并令
$\overline{V}_j$ 为 $V_j$ 在 $X$ 中的闭包。则
$$
\alpha \cap \beta = \sum m_j[\overline{V}_j]
$$
$\beta$ 的基变换是 $Y' = Y \times_S S'$ 上的循环
$\beta' = \sum [Z \times_S S']_{e + c}$。令
$Z'_a \subset Z \times_S S'$ 为不可约分支，记 $y'_a \in Z'_a$ 为其
一般点，并记 $n_a$ 为 $Z'_a$ 在 $Z \times_S S'$ 中的重数。我们有
$$
\beta' = \sum [Z \times_S S']_{e + c} = \sum n_a[Z'_a]
$$
由于 $\alpha'$ 是 $\alpha$ 沿 $Y' \to Y$ 的基变换，故
$\alpha'_{y'_a} = \sum m_j [V_{j, \kappa(y'_a)}]_r$。令
$V'_{jab} \subset V_{j, \kappa(y'_a)}$ 为不可约分支，并记 $m_{jab}$
为 $V'_{jab}$ 在 $V_{j, \kappa(y'_a)}$ 中的重数。我们有
$$
\alpha'_{y'_a} = \sum m_j [V_{j, \kappa(y'_a)}]_r =
\sum m_j m_{jab} [V'_{jab}]
$$
因此
$$
\alpha' \cap \beta' = \sum n_a m_j m_{jab} [\overline{V}'_{jab}]
$$
其中 $\overline{V}'_{jab}$ 是 $V'_{jab}$ 在 $X'$ 中的闭包。因此，为证明
所需等式，只需证明
\begin{enumerate}
\item $\overline{V}_j \times_S S'$ 的不可约分支是概形
$\overline{V}'_{jab}$，并且
\item $\overline{V}'_{jab}$ 在 $\overline{V}_j \times_S S'$ 中的重数等于
$n_a m_{jab}$。
\end{enumerate}
注意，$V_j \to \overline{V}_j$ 是整概形间的双有理态射。态射
$V_j \times_S S' \to V_j$ 与
$\overline{V}_j \times_S S' \to \overline{V}_j$ 都是平坦的，因而把
不可约分支的一般点分别映到 $V_j$ 与 $\overline{V}_j$ 的（唯一）一般点。
由此，$V_j \times_S S' \to \overline{V}_j \times_S S'$ 是双有理态射，
故在不可约分支之间诱导双射，并识别其重数。这意味着只需证明
$V_j \times_S S'$ 的不可约分支是概形 $V'_{jab}$，且 $V'_{jab}$ 在
$V_j \times_S S'$ 中的重数等于 $n_a m_{jab}$。然而，这正是在断言图
$$
\xymatrix{
Z_r(V_j) \ar[r] & Z_{r + c}(V_j \times_S S') \\
Z_0(\Spec(\kappa(y))) \ar[r] \ar[u] &
Z_c(\Spec(\kappa(y)) \times_S S') \ar[u]
}
$$
交换，其中水平箭头是沿
$\Spec(\kappa(y)) \times_S S' \to \Spec(\kappa(y))$ 的基变换，竖直箭头
是平坦拉回。这已在 Chow 同调，引理
\ref{chow-lemma-pullback-base-change-pullback} 中证明。

\medskip\noindent
引理中关于平坦拉回的陈述含义如下。令 $(S, \delta)$、$X \to Y$、
$\alpha$ 与 $\beta$ 如上面构造 $\alpha \cap \beta$ 时所述。设
$Y' \to Y$ 是局部有限型、相对维数为 $c$ 的平坦态射。令 $\alpha'$ 为
$\alpha$ 沿 $Y' \to Y$ 的基变换，令 $\beta'$ 为 $\beta$ 的平坦拉回。
与平坦拉回相容，是指 $\alpha' \cap \beta'$ 是
$\alpha \cap \beta$ 沿 $X \times_Y Y' \to Y$ 的平坦拉回。事实上，令
$S = Y$ 且 $S' = Y'$，这就是上面讨论的特例。
\end{proof}

\begin{lemma}
\label{relative-cycles-lemma-action-coherent}
令 $(S, \delta)$ 与 $f : X \to Y$ 如上。设 $\mathcal{F}$ 为凝聚
$\mathcal{O}_X$-模，且对所有 $y \in Y$ 都有
$\dim(\text{Supp}(\mathcal{F}_y)) \leq r$。设 $\mathcal{G}$ 为凝聚
$\mathcal{O}_Y$-模，且 $\dim_\delta(\text{Supp}(\mathcal{G})) \leq e$。
令 $\alpha = [\mathcal{F}/X/Y]_r$（例
\ref{relative-cycles-example-family-associated-module}），并令 $\beta = [\mathcal{G}]_e$
（Chow 同调，定义
\ref{chow-definition-cycle-associated-to-coherent-sheaf}）。若
$\mathcal{F}$ 在 $Y$ 上平坦，则 $\alpha \cap \beta =
[\mathcal{F} \otimes_{\mathcal{O}_X} f^*\mathcal{G}]_{r + e}$。
\end{lemma}

\begin{proof}
注意
$$
\text{Supp}(\mathcal{F} \otimes_{\mathcal{O}_X} f^*\mathcal{G}) =
\text{Supp}(\mathcal{F}) \cap f^{-1}\text{Supp}(\mathcal{G}) =
\bigcup\nolimits_{y \in \text{Supp}(\mathcal{G})} \text{Supp}(\mathcal{F}_y)
$$
由此，它是 $\delta$-维数 $\leq r + e$ 的闭子集。因此表达式
$[\mathcal{F} \otimes_{\mathcal{O}_X} f^*\mathcal{G}]_{r + e}$ 有意义。

\medskip\noindent
沿用构造 $\alpha \cap \beta$ 时引入的记号
$\beta = \sum n_i[Z_i]$、$y_i \in Z_i$、
$\alpha_{y_i} = \sum m_{ij} [V_{ij}]$ 与 $\overline{V}_{ij}$。由于
$\beta = [\mathcal{G}]_e$，可知 $Z_i$ 是
$\text{Supp}(\mathcal{G})$ 中 $\delta$-维数为 $e$ 的不可约分支。
类似地，$V_{ij}$ 是 $\text{Supp}(\mathcal{F}_{y_i})$ 中维数为 $r$ 的
不可约分支。结合这一点与第一段中的等式可知，$\overline{V}_{ij}$ 是
$\text{Supp}(\mathcal{F} \otimes_{\mathcal{O}_X} f^*\mathcal{G})$
中 $\delta$-维数为 $r + e$ 的不可约分支。因此，为证明本引理，只需证明
$$
\text{length}_{\mathcal{O}_{X, \xi_{ij}}}(
(\mathcal{F} \otimes_{\mathcal{O}_X} f^*\mathcal{G})_{\xi_{ij}})
=
\text{length}_{\mathcal{O}_{X_{y_i}, \xi_{ij}}}((\mathcal{F}_{y_i})_{\xi_{ij}})
\cdot
\text{length}_{\mathcal{O}_{Y, y_i}}(\mathcal{G}_{y_i})
$$
由证明第一段，左端等于下述 $B = \mathcal{O}_{X, \xi_{ij}}$-模的长度：
$$
\mathcal{G}_{y_i}
\otimes_{\mathcal{O}_{Y, y_i}}
\mathcal{F}_{\xi_{ij}} =
M \otimes_A N
$$
这里，$M = \mathcal{G}_{y_i}$ 是有限长度的
$A = \mathcal{O}_{Y, y_i}$-模，而 $N = \mathcal{F}_{\xi_{ij}}$ 是有限
$B$-模，并且 $N/\mathfrak m_AN$ 长度有限。由于 $\mathcal{F}$ 在 $Y$
上平坦，模 $N$ 是 $A$-平坦的。公式右端等于
$$
\text{length}_B(N/\mathfrak m_A N) \cdot \text{length}_A(M)
$$
因此，公式左端与右端都关于 $M$ 可加（使用 $N$ 在 $A$ 上的平坦性）。
故只需在 $M = \kappa_A$ 为剩余域时证明该公式，而此时结论立即成立。
\end{proof}

\begin{lemma}
\label{relative-cycles-lemma-action-closed}
令 $(S, \delta)$ 与 $f : X \to Y$ 如上。设 $Z \subset X$ 为在 $Y$
上相对维数 $\leq r$ 的闭子概形。令 $\alpha = [Z/X/Y]_r$（例
\ref{relative-cycles-example-family-associated-closed}）。设 $W \subset Y$ 为
$\delta$-维数 $\leq e$ 的闭子概形。令 $\beta = [W]_e$（Chow 同调，
定义 \ref{chow-definition-cycle-associated-to-closed-subscheme}）。若
$Z$ 在 $Y$ 上平坦，则 $\alpha \cap \beta = [Z \times_Y W]_{r + e}$。
\end{lemma}

\begin{proof}
在引理 \ref{relative-cycles-lemma-action-coherent} 中取
$\mathcal{F} = \mathcal{O}_Z$ 和 $\mathcal{F} = \mathcal{O}_W$ 即得。
\end{proof}

\begin{lemma}
\label{relative-cycles-lemma-flat-pullback-as-action}
令 $(S, \delta)$ 与 $f : X \to Y$ 如上。假设 $f$ 平坦，且相对维数为
$e$。对 $\beta \in Z_r(Y)$，在 $Z_{e + r}(X)$ 中有
$f^*\beta = [X/X/Y]_e \cap \beta$。
\end{lemma}

\begin{proof}
只需在 $X$ 上局部地证明该等式，故可设 $Y$ 仿射。此时 $\beta$ 是素循环的
有限整数线性组合。由于等式两端都关于 $\beta$ 可加，可设
$\beta = [W]$，其中 $W$ 是 $Y$ 中 $\delta$-维数为 $r$ 的整闭子概形。
令 $f^{-1}(W) = W \times_Y X$ 为 $W$ 在 $X$ 中的概形论逆像。由引理
\ref{relative-cycles-lemma-action-closed}，有
$[X/X/Y]_e \cap \beta = [W \times_Y X]_{r + e}$；而由 Chow 同调，定义
\ref{chow-definition-flat-pullback}，有
$f^*\beta = [f^{-1}(W)]_{r + e}$。故得到所需等式。
\end{proof}

\begin{lemma}
\label{relative-cycles-lemma-action-push-pull}
令 $(S, \delta)$ 如上。设
$$
\xymatrix{
X' \ar[r]_f \ar[d] & X \ar[d] \\
Y' \ar[r]^g & Y
}
$$
是 $S$ 上局部有限型概形的笛卡尔图，且 $g$ 固有。设 $r, e \geq 0$，
令 $\alpha$ 为 $X/Y$ 的纤维上的 $r$-循环族，并令
$\beta' \in Z_e(Y')$。则
$f_*(g^*\alpha \cap \beta') = \alpha \cap g_*\beta'$。
\end{lemma}

\begin{proof}
由于要证明的是 $X$ 上循环的等式，由引理 \ref{relative-cycles-lemma-action-opens}，可在
$Y$ 上局部地工作。因此可设 $Y$ 仿射，从而 $Y'$ 拟紧。特别地，
$\beta'$ 是素循环的有限线性组合。由于 $- \cap -$ 关于第二个变量线性
（引理 \ref{relative-cycles-lemma-action-bilinear}），只需在 $\beta' = [Z']$ 时证明该
等式，其中 $Z' \subset Y'$ 是 $\delta$-维数为 $e$ 的整闭子概形。令
$Z = g(Z')$。这是 $Y$ 中 $\delta$-维数 $\leq e$ 的整闭子概形。为简便
起见，我们假设 $Z$ 的 $\delta$-维数等于 $e$，并把另一种（更容易的）
情形留给读者。令 $y \in Z$ 与 $y' \in Z'$ 为一般点。写
$\alpha_y = \sum m_j[V_j]$，其中 $V_j \subset X_y$ 是维数为 $r$ 的
整闭子概形。

\medskip\noindent
先假设 $g$ 为闭浸入。于是 $g_*\beta' = [Z]$ 且
$(g^*\alpha)_{y'} = \sum n_j[V_j]$；这是有意义的，因为 $V_j$ 包含在
$X_y$ 的闭子概形 $X'_{y'}$ 中。因此，此时等式显然：两边都得到
$\sum m_j[\overline{V}_j]$，其中 $\overline{V}_j$ 是 $V_j$ 在闭子概形
$X' \subset X$ 中的闭包。

\medskip\noindent
回到上述 $\beta' = [Z']$ 的一般情形。令 $W = Z \times_Y X$ 且
$W' = Z' \times_{Y'} X'$。考虑笛卡尔方块
$$
\xymatrix{
W \ar[r] \ar[d] & X \ar[d] \\
Z \ar[r] & Y
}
\quad
\xymatrix{
W' \ar[r] \ar[d] & X' \ar[d] \\
Z' \ar[r] & Y'
}
\quad
\xymatrix{
W' \ar[r] \ar[d] & W \ar[d] \\
Z' \ar[r] & Z
}
$$
由上一段已知前两个方块的结论，一个形式论证表明，只需对最后一个方块与
元素 $\beta' = [Z'] \in Z_e(Z')$ 证明结论。于是归结为下一段讨论的情形。

\medskip\noindent
假设 $Y' \to Y$ 是 $\delta$-维数为 $e$ 的整概形之间的一般有限态射，且
$\beta' = [Y']$。此时，$f_*(g^*\alpha \cap \beta')$ 与
$\alpha \cap g_*\beta'$ 都可写成支配 $Y$ 的素循环之和。因此，为检验
等式，可以把 $Y$ 换成一个非空开子概形。作此替换后，可设 $g$ 有限且
平坦，譬如次数为 $d \geq 1$。当然，这意味着
$g_*\beta' = g_*[Y'] = d[Y]$。又有 $\beta' = [Y'] = g^*[Y]$。因此
$$
f_*(g^*\alpha \cap \beta') =
f_*(g^*\alpha \cap g^*[Y]) =
f_*f^*(\alpha \cap [Y]) =
d (\alpha \cap [Y]) =
\alpha \cap g_*\beta'
$$
正合所需。第二个等号来自引理 \ref{relative-cycles-lemma-action-base-change}，第三个等号
来自 Chow 同调，引理 \ref{chow-lemma-finite-flat}。
\end{proof}







\section{Chow 群上的作用}
\label{relative-cycles-section-action-chow}

\noindent
当 $\alpha$ 为相对 $r$-循环时，第 \ref{relative-cycles-section-action} 节的运算
$\alpha \cap -$ 可经由有理等价因子化，并定义一个双变类。

\begin{lemma}
\label{relative-cycles-lemma-closed-in-X-gysin}
令 $(S, \delta)$ 如第 \ref{relative-cycles-section-action} 节所述。设
$f : X' \to X$ 为 $S$ 上局部有限型概形之间的固有态射。令
$(\mathcal{L}, s, i : D \to X)$ 如 Chow 同调，定义
\ref{chow-definition-gysin-homomorphism} 所述。构造图
$$
\xymatrix{
D' \ar[d]_g \ar[r]_{i'} & X' \ar[d]^f \\
D \ar[r]^i & X
}
$$
如 Chow 同调，注记 \ref{chow-remark-pullback-pairs} 所述。若
$\mathcal{L}|_D \cong \mathcal{O}_D$，则对任意
$\alpha' \in Z_{k + 1}(X')$，在 $Z_k(D)$ 中有
$i^*f_*\alpha' = g_*(i')^*\alpha'$。
\end{lemma}

\begin{proof}
该陈述有意义，因为所有运算都在循环层面定义；关于 Gysin 映射，参见
Chow 同调，注记 \ref{chow-remark-gysin-on-cycles}。假设对某个整闭子概形
$W' \subset X'$ 有 $\alpha = [W']$。令 $W = f(W') \subset X$。若
$W' \not \subset D'$，则 $W \not \subset D$，并且
$$
[W' \cap D']_k = \text{div}_{\mathcal{L}'|_{W'}}({s'|_{W'}})
\quad\text{且}\quad
[W \cap D]_k = \text{div}_{\mathcal{L}|_W}(s|_W)
$$
因此，由 Chow 同调，引理 \ref{chow-lemma-equal-c1-as-cycles}，第一个循环
在 $f_*$ 下的像等于第二个循环。故等式在循环层面成立。若
$W' \subset D'$，则 $W \subset D$，且按构造两端都为零。
\end{proof}

\begin{lemma}
\label{relative-cycles-lemma-action-gysin}
令 $(S, \delta)$ 如第 \ref{relative-cycles-section-action} 节所述。设 $X \to Y$ 为
$S$ 上局部有限型概形之间的态射。设 $r \geq 0$，并令
$\alpha \in z(X/Y, r)$ 为 $X/Y$ 上的相对 $r$-循环。令
$(\mathcal{L}, s, i : D \to Y)$ 如 Chow 同调，定义
\ref{chow-definition-gysin-homomorphism} 所述。构造笛卡尔图
$$
\xymatrix{
E \ar[d] \ar[r]_j & X \ar[d] \\
D \ar[r]^i & Y
}
$$
参见 Chow 同调，注记 \ref{chow-remark-pullback-pairs}。若
$\mathcal{L}|_D \cong \mathcal{O}_D$，则对 $e \in \mathbf{Z}$，图
$$
\xymatrix{
Z_e(D) \ar[rr]_{i^*\alpha \cap -} & &
Z_{e + r}(E) \\
Z_{e + 1}(Y) \ar[u]^{i^*} \ar[rr]^{\alpha \cap -} & &
Z_{r + e + 1}(X) \ar[u]_{j^*}
}
$$
交换，其中竖直箭头 $i^*$ 与 $j^*$ 是 Chow 同调，注记
\ref{chow-remark-gysin-on-cycles} 中的循环 Gysin 映射。
\end{lemma}

\begin{proof}
先作预备说明。假设 $g : Y' \to Y$ 为包络（Chow 同调，定义
\ref{chow-definition-envelope}）。分别用 $D', i', E', j', X', \alpha'$
表示 $D, i, E, j, X, \alpha$ 沿 $g$ 的基变换，并记
$f : X' \to X$ 为投影。假设本引理对
$D', i', E', j', X', Y', \alpha'$ 成立。于是，若
$\beta' \in Z_{e + 1}(Y')$，则
\begin{align*}
i^*\alpha \cap i^*g_*\beta'
& =
i^*\alpha \cap f_*(i')^*\beta' \\
& =
f_*(f^*i^*\alpha \cap (i')^*\beta') \\
& =
f_*((i')^*\alpha' \cap (i')^*\beta') \\
& =
f_*((j')^*(\alpha' \cap \beta')) \\
& =
j^*(f_*(f^*\alpha \cap \beta')) \\
& =
j^*(\alpha \cap g_*\beta')
\end{align*}
这里，第一个等号来自引理 \ref{relative-cycles-lemma-closed-in-X-gysin}，第二个等号来自
引理 \ref{relative-cycles-lemma-action-push-pull}，第三个等号是 $\alpha'$ 的定义，第四个
等号来自本引理对 $D', i', E', j', X', \alpha'$ 成立的假设，第五个等号
来自引理 \ref{relative-cycles-lemma-closed-in-X-gysin}，第六个等号来自引理
\ref{relative-cycles-lemma-action-push-pull}。因此，本引理对
$g_* : Z_{e + 1}(Y') \to Z_e(Y)$ 的像成立。然而，由于 $g$ 完全分解，
此映射满射，故本引理对 $D, i, E, j, X, Y, \alpha$ 成立。

\medskip\noindent
令 $\beta \in Z_{e + 1}(Y)$。我们须证明，在 $E$ 上作为循环有
$(D \to Y)^*\alpha \cap i^*\beta = j^*(\alpha \cap \beta)$。这一问题在
$E$ 上是局部的，因而可以把 $X$ 与 $Y$ 换成开子概形。（这里用到运算
$i^*$、$j^*$、$\alpha \cap - $ 与 $(D \to Y)^*\alpha \cap -$ 的构造
均与局部化交换。对 Gysin 映射这是显然的；对其余运算则由引理
\ref{relative-cycles-lemma-action-opens} 得到。）因此可设 $X$ 与 $Y$ 仿射，并归结到
下一段讨论的情形。

\medskip\noindent
假设 $X$ 与 $Y$ 拟紧。由证明第一段与引理 \ref{relative-cycles-lemma-get-cycles}，还可设
$\alpha$ 属于 (\ref{relative-cycles-equation-cycle-classes}) 的像。由所涉运算的线性，
可设 $\alpha = [Z/X/Y]_r$，其中 $Z \subset X$ 是在 $Y$ 上平坦且相对维数
$\leq r$ 的闭子概形。又因 $Y$ 拟紧，循环 $\beta$ 是素循环的有限线性
组合。由于所涉运算是线性的，只需在 $\beta = [W]$ 时证明等式，其中
$W \subset Y$ 是 $\delta$-维数为 $e + 1$ 的整闭子概形。

\medskip\noindent
若 $W \subset D$，则一方面 $i^*[W] = 0$，另一方面
$\alpha \cap [W]$ 支撑在 $E$ 上，故也有
$j^*(\alpha \cap [W]) = 0$。因此此时等式成立。

\medskip\noindent
设 $W \not \subset D$。于是 $i^*[W] = [D \cap W]_e$。注意，
$\alpha = [Z/X/Y]_r$ 沿 $i$ 的拉回 $i^*\alpha$ 是
$[(E \cap Z)/E/D]_r$，而且
$(E \cap Z) = E \times_Y Z = D \times_Y Z$ 在 $D$ 上平坦。因此，两次
使用引理 \ref{relative-cycles-lemma-action-closed} 可得
$$
i^*\alpha \cap i^*[W] =
[(E \cap Z) \times_D (D \cap W)]_{r + e} =
[E \cap (Z \times_Y W)]_{r + e} =
j^*(\alpha \cap [W])
$$
正合所需。
\end{proof}

\begin{proposition}
\label{relative-cycles-proposition-get-bivariant-class}
令 $(S, \delta)$ 如第 \ref{relative-cycles-section-action} 节所述。设 $X \to Y$ 为
$S$ 上局部有限型概形之间的态射。设 $r \geq 0$，并令
$\alpha \in z(X/Y, r)$ 为 $X/Y$ 上的相对 $r$-循环。对每个局部有限型
态射 $g : Y' \to Y$ 与每个 $e \in \mathbf{Z}$，考虑运算
$$
g^*\alpha \cap - : Z_e(Y') \to Z_{r + e}(X')
$$
其中 $X' = Y' \times_Y X$。这族运算可经由有理等价因子化，从而定义双变类
$c(\alpha) \in A^{-r}(X \to Y)$。
\end{proposition}

\begin{proof}
由引理 \ref{relative-cycles-lemma-action-gysin} 与 Chow 同调，引理
\ref{chow-lemma-factors-through-rational-equivalence}，该运算可经由有理等价
因子化。由 Chow 同调，引理 \ref{chow-lemma-bivariant-weaker} 以及引理
\ref{relative-cycles-lemma-action-push-pull}、\ref{relative-cycles-lemma-action-base-change} 和
\ref{relative-cycles-lemma-action-gysin}，所得 Chow 群上的运算是双变类。
\end{proof}

\begin{remark}
\label{relative-cycles-remark-characterize-relative-cycles}
令 $(S, \delta)$ 如第 \ref{relative-cycles-section-action} 节所述。设 $X \to Y$ 为
$S$ 上局部有限型概形之间的态射，并设 $r \geq 0$。设 $c$ 是一条规则，
它给每个局部有限型态射 $g : Y' \to Y$ 与每个 $e \in \mathbf{Z}$ 配置运算
$$
c \cap - : Z_e(Y') \to Z_{r + e}(X')
$$
并且它按引理 \ref{relative-cycles-lemma-action-gysin} 的意义与固有前推、平坦拉回和
Gysin 映射相容。则我们断言，存在 $X/Y$ 上的相对 $r$-循环 $\alpha$，
使得对上述每个 $g$ 都有 $c \cap = g^*\alpha \cap -$。若以后需要这一
结论，我们将在此仔细陈述并证明它。
\end{remark}









\section{纤维上循环族的复合}
\label{relative-cycles-section-compose-families}

\noindent
设 $X \to Y \to S$ 为概形态射，且二者均局部有限型。设
$r, e \geq 0$。令 $\alpha$ 为 $X/Y$ 的纤维上的 $r$-循环族，令
$\beta$ 为 $Y/S$ 的纤维上的 $e$-循环族。定义
$$
(\alpha \circ \beta)_s = (Y_s \to Y)^*\alpha \cap \beta_s
$$
便得到 $X/S$ 的纤维上的 $(r + e)$-循环族 $\alpha \circ \beta$。更准确地，
表达式 $(Y_s \to Y)^*\alpha$ 表示 $\alpha$ 沿 $Y_s \to Y$ 的基变换，
所得为 $X_s/Y_s$ 的纤维上的 $r$-循环族；运算 $- \cap -$ 已在第
\ref{relative-cycles-section-action} 节中定义并研究\footnote{明确地说，我们把
$s = \Spec(\kappa(s))$ 用作底概形，并取 $\delta(s) = 0$。}。

\begin{lemma}
\label{relative-cycles-lemma-construction-bilinear}
上述构造是双线性的，即有
$(\alpha_1 + \alpha_2) \circ \beta \alpha_1 \circ \beta +
\alpha_1 \circ \beta$，且 $\alpha \circ (\beta_1 + \beta_2) =
\alpha \circ \beta_1 + \alpha \circ \beta_2$.
\end{lemma}

\begin{proof}
略。提示：由引理 \ref{relative-cycles-lemma-action-bilinear}，该构造在纤维上双线性。
\end{proof}

\begin{lemma}
\label{relative-cycles-lemma-construction-opens}
若 $U \subset X$ 与 $V \subset Y$ 为开集，且 $f(U) \subset V$，则
$(\alpha \circ \beta)|_U$ 等于 $\alpha|_U \circ \beta|_V$。
\end{lemma}

\begin{proof}
略。提示：在纤维上使用引理 \ref{relative-cycles-lemma-action-opens}。
\end{proof}

\begin{lemma}
\label{relative-cycles-lemma-construction-base-change}
$\alpha \circ \beta$ 的构造与基变换相容。
\end{lemma}

\begin{proof}
设 $g : S' \to S$ 为概形态射。记 $X' \to Y'$ 为 $X \to Y$ 沿 $g$ 的
基变换，记 $\alpha'$ 为 $\alpha$ 沿 $Y' \to Y$ 的基变换，并记
$\beta'$ 为 $\beta$ 沿 $S' \to S$ 的基变换。断言的含义是，
$\alpha' \circ \beta'$ 是 $\alpha \circ \beta$ 沿
$g : S' \to S$ 的基变换。

\medskip\noindent
设 $s' \in S'$ 为像是 $s \in S$ 的点。则
$$
(\alpha' \circ \beta')_{s'} = (Y'_{s'} \to Y')^*\alpha' \cap \beta'_{s'}
$$
注意
$$
(Y'_{s'} \to Y')^*\alpha' =
(Y'_{s'} \to Y')^*(Y' \to Y)^*\alpha =
(Y'_{s'} \to Y_s)^*(Y_s \to Y)^*\alpha
$$
而 $\beta'_{s'}$ 是 $\beta_s$ 沿
$s' = \Spec(\kappa(s')) \to \Spec(\kappa(s)) = s$ 的基变换。因此，把引理
\ref{relative-cycles-lemma-action-base-change} 应用于 $(Y_s \to Y)^*\alpha$、$\beta_s$、
$X_s \to Y_s \to s$ 以及沿 $s' \to s$ 的基变换，即得结论。
\end{proof}

\begin{lemma}
\label{relative-cycles-lemma-construction-coherent}
设 $f : X \to Y$ 与 $Y \to S$ 为概形态射，且二者均局部有限型。设
$r, e \geq 0$。设 $\mathcal{F}$ 为有限型拟凝聚 $\mathcal{O}_X$-模，且对
所有 $y \in Y$ 有 $\dim(\text{Supp}(\mathcal{F}_y)) \leq r$。设
$\mathcal{G}$ 为有限型拟凝聚 $\mathcal{O}_Y$-模，且对所有 $s \in S$ 有
$\dim(\text{Supp}(\mathcal{G}_s)) \leq e$。若
$\alpha = [\mathcal{F}/X/Y]_r$ 且 $\beta = [\mathcal{G}/Y/S]_e$（例
\ref{relative-cycles-example-family-associated-module}），并且 $\mathcal{F}$ 在 $Y$ 上
平坦，则 $\alpha \circ \beta =
[\mathcal{F} \otimes_{\mathcal{O}_X} f^*\mathcal{G}/X/S]_{r + e}$。
\end{lemma}

\begin{proof}
首先，$\mathcal{F} \otimes_{\mathcal{O}_X} f^*\mathcal{G}$ 是有限型拟凝聚
$\mathcal{O}_X$-模。设 $s \in S$。注意
$$
(\mathcal{F} \otimes_{\mathcal{O}_X} f^*\mathcal{G})_s =
\mathcal{F}_s \otimes_{\mathcal{O}_{X_s}} f_s^*\mathcal{G}_s
$$
这来自张量积的右正合性。此外，$\mathcal{F}_s$ 作为平坦模的基变换，在
$Y_s$ 上平坦。因此，由引理 \ref{relative-cycles-lemma-action-coherent} 得到等式
$(\alpha \circ \beta)_s =
[(\mathcal{F} \otimes_{\mathcal{O}_X} f^*\mathcal{G})_s]_{r + e}$。
\end{proof}

\begin{lemma}
\label{relative-cycles-lemma-construction-closed}
设 $f : X \to Y$ 与 $Y \to S$ 为概形态射，且二者均局部有限型。设
$r, e \geq 0$。设 $Z \subset X$ 为在 $Y$ 上相对维数 $\leq r$ 的闭子概形，
并设 $W \subset Y$ 为在 $S$ 上相对维数 $\leq e$ 的闭子概形。若
$\alpha = [Z/X/Y]_r$ 且 $\beta = [W/Y/S]_e$（例
\ref{relative-cycles-example-family-associated-closed}），并且 $Z$ 在 $Y$ 上平坦，则
$\alpha \circ \beta = [Z \times_Y W/X/S]_{r + e}$。
\end{lemma}

\begin{proof}
在引理 \ref{relative-cycles-lemma-construction-coherent} 中取
$\mathcal{F} = \mathcal{O}_Z$ 与 $\mathcal{F} = \mathcal{O}_W$ 即得。
\end{proof}

\begin{lemma}
\label{relative-cycles-lemma-flat-pullback-as-composition}
设 $f : X \to Y$ 与 $Y \to S$ 为概形态射，且二者均局部有限型。假设
$f$ 平坦，且相对维数为 $e$。令 $\beta$ 为 $Y/S$ 的纤维上的
$r$-循环族。则作为 $X/S$ 的纤维上的 $(e + r)$-循环族，有
$f^*\beta = [X/X/Y]_e \circ \beta$。
\end{lemma}

\begin{proof}
展开定义，并使用 $[X/X/Y]_e$ 的构造与基变换相容这一事实（引理
\ref{relative-cycles-lemma-family-associated-closed}），即可由引理
\ref{relative-cycles-lemma-flat-pullback-as-action} 得到结论。
\end{proof}

\begin{lemma}
\label{relative-cycles-lemma-construction-push-pull}
设 $S$ 为概形。设
$$
\xymatrix{
X' \ar[r]_f \ar[d] & X \ar[d] \\
Y' \ar[r]^g & Y
}
$$
是 $S$ 上局部有限型概形的笛卡尔图，且 $g$ 固有。设 $r, e \geq 0$。
令 $\alpha$ 为 $X/Y$ 的纤维上的 $r$-循环族，并令 $\beta'$ 为 $Y'/S$
的纤维上的 $e$-循环族。则
$f_*(g^*(\alpha) \circ \beta') = \alpha \circ g_*\beta'$。
\end{lemma}

\begin{proof}
展开定义后，结论来自引理 \ref{relative-cycles-lemma-action-push-pull}。
\end{proof}

\begin{lemma}
\label{relative-cycles-lemma-construction-composition}
令 $(S, \delta)$ 如 Chow 同调，情形 \ref{chow-situation-setup} 所述。设
$X \to Y \to Z$ 为 $S$ 上局部有限型概形之间的态射。设
$r, s, e \geq 0$。则
$$
(\alpha \circ \beta) \cap \gamma = \alpha \cap (\beta \cap \gamma)
\quad\text{在}\quad Z_{r + s + e}(X)
$$
其中 $\alpha$ 为 $X/Y$ 的纤维上的 $r$-循环族，$\beta$ 为 $Y/Z$ 的纤维
上的 $s$-循环族，且 $\gamma \in Z_e(Z)$。
\end{lemma}

\begin{proof}
由于要证明的是 $X$ 上循环的等式，由引理 \ref{relative-cycles-lemma-action-opens}，可在
$Z$ 上局部地工作。因此可设 $Z$ 仿射。特别地，$\gamma$ 是素循环的有限
线性组合。由于 $- \cap -$ 关于第二个变量线性（引理
\ref{relative-cycles-lemma-action-bilinear}），只需在 $\gamma = [W]$ 时证明等式，其中
$W \subset Z$ 是 $\delta$-维数为 $e$ 的整闭子概形。

\medskip\noindent
令 $z \in W$ 为一般点。在 $Z_s(Y_z)$ 中写
$\beta_z = \sum m_j[V_j]$。于是 $\beta \cap \gamma$ 等于
$\sum m_j[\overline{V}_j]$，其中 $\overline{V}_j \subset Y$ 是一个在
$Y \to Z$ 下映入 $W$、且一般纤维为 $V_j$ 的整闭子概形。令
$y_j \in V_j$ 为一般点。我们也把它看成 $\overline{V}_j$ 的一般点
（它映到 $W$ 中的 $z$）。在 $Z_r(X_{y_j})$ 中写
$\alpha_{y_j} = \sum n_{jk} [W_{jk}]$。则
$\alpha \cap (\beta \cap \gamma)$ 等于
$$
\sum m_j n_{jk} [\overline{W}_{jk}]
$$
其中 $\overline{W}_{jk} \subset X$ 是一个在 $X \to Y$ 下映入
$\overline{V}_j$、且一般纤维为 $W_{jk}$ 的整闭子概形。

\medskip\noindent
另一方面，考虑
$$
(\alpha \circ \beta)_z = (Y_z \to Y)^*\alpha \cap \beta_z =
(Y_z \to Y)^*\alpha \cap (\sum m_j [V_j])
$$
由 $- \cap -$ 的构造，它等于 $X_z$ 上的循环
$$
\sum m_j n_{jk} [(\overline{W}_{jk})_z]
$$
因此按定义得到
$$
(\alpha \circ \beta) \cap [W] =
\sum m_j n_{jk} [\widetilde{W}_{jk}]
$$
其中 $\widetilde{W}_{jk} \subset X$ 是一个在 $X \to Z$ 下映入 $W$、且
一般纤维为 $(\overline{W}_{jk})_z$ 的整闭子概形。显然必有
$\widetilde{W}_{jk} = \overline{W}_{jk}$，证明完成。
\end{proof}









\section{相对循环的复合}
\label{relative-cycles-section-compose}

\noindent
设 $S$ 为局部 Noether 概形。设 $X \to Y$ 为 $S$ 上局部有限型概形之间的
态射。我们将利用第 \ref{relative-cycles-section-compose-families} 节的构造定义映射
$$
z(X/Y, r) \otimes_\mathbf{Z} z(Y/S, e) \longrightarrow z(X/S, r + e),\quad
\alpha \otimes \beta \longmapsto \alpha \circ \beta
$$
我们已经知道该构造是双线性的（引理 \ref{relative-cycles-lemma-construction-bilinear}），
故只要证明下述结论，就得到所显示的箭头。

\begin{lemma}
\label{relative-cycles-lemma-well-defined}
若 $\alpha$ 与 $\beta$ 都是相对循环，则 $\alpha \circ \beta$ 也是相对循环。
\end{lemma}

\begin{proof}
由引理 \ref{relative-cycles-lemma-construction-base-change}，$\alpha \circ \beta$ 的构造与
基变换相容。因此可设 $S$ 是某个离散赋值环的谱，其一般点为 $\eta$、
闭点为 $0$；我们须证明
$sp_{X/S}((\alpha \circ \beta)_\eta) = (\alpha \circ \beta)_0$。由于要
证明的是循环的等式，可在 $Y$ 与 $X$ 上局部地工作（这里使用引理
\ref{relative-cycles-lemma-construction-opens} 与 \ref{relative-cycles-lemma-specialization-flat-pullback}
看出各构造与限制交换）。故可设 $X$ 与 $Y$ 仿射。由引理
\ref{relative-cycles-lemma-get-cycles}，可找到固有完全分解态射 $g : Y' \to Y$，使得
$g^*\alpha$ 属于 (\ref{relative-cycles-equation-cycle-classes}) 的像。

\medskip\noindent
由于态射族 $g_\eta : Y'_\eta \to Y_\eta$ 完全分解，可找到
$\beta'_\eta \in Z_e(Y'_\eta)$，使得
$\beta_\eta = \sum g_{\eta, *}\beta'_\eta$，参见 Chow 同调，引理
\ref{chow-lemma-envelope}。令 $\beta'_0 = sp_{Y'/S}(\beta'_\eta)$，从而
$\beta' = (\beta'_\eta, \beta'_0)$ 是 $Y'/S$ 上的相对 $e$-循环。于是
$g_*\beta'$ 与 $\beta$ 是 $Y/S$ 上的相对 $e$-循环（引理
\ref{relative-cycles-lemma-relative-cycle-functoriality}），并且它们在 $\eta$ 处取值相同，
故二者相等（引理 \ref{relative-cycles-lemma-uniqueness-extension}）。由线性（引理
\ref{relative-cycles-lemma-construction-bilinear}），只需证明
$\alpha \circ g_*\beta'$ 是相对 $(r + e)$-循环。

\medskip\noindent
令 $X' = X \times_Y Y'$，并记 $f : X' \to X$ 为投影。由引理
\ref{relative-cycles-lemma-construction-push-pull}，有
$\alpha \circ g_*\beta' = f_*(g^*\alpha \circ \beta')$。由引理
\ref{relative-cycles-lemma-relative-cycle-functoriality}，只需证明
$g^*\alpha \circ \beta'$ 是相对 $(r + e)$-循环。使用引理
\ref{relative-cycles-lemma-get-cycles-dvr} 与双线性，即归结到下一段讨论的情形。

\medskip\noindent
假设 $\alpha = [Z/X/Y]_r$ 且 $\beta = [W/Y/S]$，其中 $Z \subset X$
是在 $Y$ 上平坦且相对维数 $\leq r$ 的闭子概形，而 $W \subset Y$ 是在
$S$ 上平坦且相对维数 $\leq e$ 的闭子概形。由引理
\ref{relative-cycles-lemma-construction-closed} 可知
$$
\alpha \circ \beta = [Z \times_X W/X/S]_{r + e}
$$
而 $Z \times_X W \subset X$ 是在 $S$ 上平坦、相对维数
$\leq r + e$ 的闭子概形。由引理
\ref{relative-cycles-lemma-family-associated-closed-specialization}，这是相对
$(r + e)$-循环。
\end{proof}

\begin{lemma}
\label{relative-cycles-lemma-composition-fundamental-cycles}
设 $f : X \to Y$ 与 $g : Y \to S$ 为概形态射。假设 $S$ 局部 Noether，
$g$ 局部有限型且平坦、相对维数为 $e \ge 0$，并且 $f$ 局部有限型且平坦、
相对维数为 $r \geq 0$。则在 $z(X/S, r + e)$ 中有
$[X/X/Y]_r \circ [Y/Y/S]_e = [X/X/S]_{r + e}$。
\end{lemma}

\begin{proof}
这是引理 \ref{relative-cycles-lemma-construction-closed} 的特例。
\end{proof}






\section{与 Suslin--Voevodsky 的比较}
\label{relative-cycles-section-compare}

\noindent
除循环记号取自 Chow 同调，第 \ref{chow-section-cycles} 节及后续内容外，
我们尽量采用与 \cite{SV} 相同的记号。比较如下：
\begin{enumerate}
\item \cite[Section 3.1]{SV} 中有“相对循环”、“维数为 $r$ 的相对循环”
以及“维数为 $r$ 的等维相对循环”的概念。本章中没有相应概念。因此，各群
$Cycl(X/S, r)$, $Cycl_{equi}(X/S, r)$,
$PropCycl(X/S, r)$ 与 $PropCycl_{equi}(X/S, r)$
在本章中没有对应物。
\item 在 \cite[page 36]{SV} 页底定义了各群
$z(X/S, r)$, $c(X/S, r)$, $z_{equi}(X/S, r)$, $c_{equi}(X/S, r)$
。当 $S$ 是分离 Noether 概形且 $X \to S$ 分离有限型时，这些定义与我们的
概念一致。
\item 在 \cite{SV} 中，符号 $z(X/S, r)$ 有时用来表示 $S$ 上有限型概形
范畴上的预层 $S' \mapsto z(S' \times_S X/S', r)$。符号
$c(X/S, r)$、$z_{equi}(X/S, r)$ 与 $c_{equi}(X/S, r)$。
也作类似使用。
\item \cite{SV} 中定义的基变换、平坦拉回和固有前推，在两套定义都适用时
与我们的定义一致。
\item 对 $\alpha \in z(X/S, r)$，第 \ref{relative-cycles-section-action} 节定义的运算
$\alpha \cap - : Z_e(S) \to Z_{e + r}(X)$，在两者都有定义时，与
\cite[Section 3.7]{SV} 中的运算 $Cor(\alpha, -)$ 一致。
\item 对 $X \to Y \to S$，第 \ref{relative-cycles-section-compose} 节定义的复合律
$z(X/Y, r) \otimes_\mathbf{Z} z(Y/S, e) \longrightarrow z(X/S, r + e)$
与 \cite[Corollary 3.7.5]{SV} 中的运算 $Cor_{X/Y}(-, -)$ 一致。
\end{enumerate}











\section{非 Noether 情形下的相对循环}
\label{relative-cycles-section-non-noetherian}

\noindent
我们强烈建议读者跳过本节。

\medskip\noindent
设 $f : X \to S$ 为有限表示的概形态射。设 $r \geq 0$。用
$Hilb(X/S, r)$ 表示所有满足下述条件的闭子概形 $Z \subset X$ 所成的集合：
$Z \to S$ 平坦、有限表示且相对维数 $\leq r$。考虑群同态
\begin{equation}
\label{relative-cycles-equation-cycle-classes-general}
\begin{matrix}
\text{自由 Abel 群} \\
\text{以 }Hilb(X/S, r)
\end{matrix}
\longrightarrow
\begin{matrix}
\text{纤维上的 }r\text{-循环族}\\
\text{对应于 }X/S
\end{matrix}
\end{equation}
它把 $\sum n_i[Z_i]$ 映到 $\sum n_i[Z_i/X/S]_r$。

\begin{lemma}
\label{relative-cycles-lemma-relative-r-cycle-general}
设 $S$ 为拟紧拟分离概形，$f : X \to S$ 为有限表示态射。设
$r \geq 0$，并令 $\alpha$ 为 $X/S$ 的纤维上的 $r$-循环族。下列条件等价：
\begin{enumerate}
\item 存在笛卡尔图
$$
\xymatrix{
X \ar[r] \ar[d] & X_0 \ar[d] \\
S \ar[r] & S_0
}
$$
其中 $X_0 \to S_0$ 是 Noether 概形之间的有限型态射，且存在
$\alpha_0 \in z(X_0/S_0, r)$，使得 $\alpha$ 是 $\alpha_0$ 沿
$S \to S_0$ 的基变换；
\item 存在有限表示的固有完全分解态射 $g : S' \to S$，使得
$g^*\alpha$ 属于 (\ref{relative-cycles-equation-cycle-classes-general}) 的像。
\end{enumerate}
\end{lemma}

\begin{proof}
给定 (1) 中的图与 $\alpha_0 \in z(X_0/S_0, r)$。由引理
\ref{relative-cycles-lemma-get-cycles}，存在固有完全分解态射 $g_0 : S'_0 \to S_0$，
使得 $g_0^*\alpha_0$ 属于 (\ref{relative-cycles-equation-cycle-classes-general}) 的像。
事实上，由于 $S'_0$ 是 Noether 的，$S'_0 \times_{S_0} X_0$ 的每个闭子概形
在 $S'_0$ 上都有限表示。令 $S' = S \times_{S_0} S'_0$，并沿
$S' \to S'_0$ 作基变换，即知 (2) 成立。

\medskip\noindent
反之，假设 (2) 成立。选取有限表示的固有完全分解态射 $g : S' \to S$，
使得 $g^*\alpha$ 属于 (\ref{relative-cycles-equation-cycle-classes-general}) 的像。令
$X' = S' \times_S X$。对某些闭子概形 $Z_a \subset X'$，写
$g^*\alpha = \sum n_a [Z_a/X'/S']_r$，其中这些子概形在 $S'$ 上平坦、
有限表示且相对维数 $\leq r$。

\medskip\noindent
把 $S = \lim S_i$ 写成转移态射为仿射态射的有向极限，其中 $S_i$ 在
$\mathbf{Z}$ 上有限型，参见极限，命题 \ref{limits-proposition-approximate}。
可找到充分大的 $i$，使得存在
\begin{enumerate}
\item 固有完全分解态射 $g_i : S'_i \to S_i$，它到 $S$ 的基变换是
$g : S' \to S$；
\item 令 $X'_i = S'_i \times_{S_i} X_i$，存在闭子概形
$Z_{ai} \subset X'_i$，其在 $S'_i$ 上平坦且相对维数 $\leq r$，并且到
$S'$ 的基变换是 $Z_a$。
\end{enumerate}
为此要使用极限中的各引理
\ref{limits-lemma-descend-finite-presentation},
\ref{limits-lemma-descend-closed-immersion-finite-presentation},
\ref{limits-lemma-descend-flat-finite-presentation},
\ref{limits-lemma-descend-surjective},
\ref{limits-lemma-eventually-proper} 和
\ref{limits-lemma-limit-dimension}
以及更多态射，引理 \ref{more-morphisms-lemma-descend-cd}。考虑
$\alpha'_i = \sum n_a [Z_{ai}/X'_i/S_i]_r \in z(X'_i/S'_i, r)$。按构造，
$\alpha'_i$ 基变换为 $X'/S'$ 的纤维上的 $r$-循环族后，与基变换
$g^*\alpha$ 一致。

\medskip\noindent
令 $S''_i = S'_i \times_{S_i} S'_i$、$X''_i = S''_i \times_{S_i} X_i$，
并令 $S'' = S' \times_S S'$、$X'' = S'' \times_S X$。用
$\text{pr}_1, \text{pr}_2 : S'' \to S'$ 与
$\text{pr}_1, \text{pr}_2 : S''_i \to S'_i$ 表示投影。$X''_i/S''_i$
上的相对 $r$-循环 $\text{pr}_1^*\alpha'_i$ 与
$\text{pr}_1^*\alpha'_i$ 基变换后，成为 $X''/S''$ 的纤维上同一个
$r$-循环族，因为
$\text{pr}_1^*g^*\alpha = \text{pr}_1^*g^*\alpha$。因此态射
$S'' \to S''_i$ 映入
$E = \{s \in S''_i : (\text{pr}_1^*\alpha'_i)_s =
(\text{pr}_1^*\alpha'_i)_s\}$。由引理 \ref{relative-cycles-lemma-relative-cycles-equal}，
这是闭子集。由于 $S'' = \lim_{i' \geq i} S''_{i'}$，由极限，引理
\ref{limits-lemma-limit-contained-in-constructible} 可知，对某个
$i' \geq i$，态射 $S''_{i'} \to S''_i$ 映入 $E$。因此，把 $i$ 换成
$i'$ 后，可设 $\text{pr}_1^*\alpha'_i = \text{pr}_1^*\alpha'_i$。由引理
\ref{relative-cycles-lemma-descend-family}，得到 $X_i/S_i$ 的纤维上唯一的 $r$-循环族
$\alpha_i$，满足 $g_i^*\alpha_i = \alpha'_i$（这里使用了
$S'_i \to S_i$ 完全分解）。由引理 \ref{relative-cycles-lemma-relative-cycles-h-descent}，
有 $\alpha_i \in z(X_i/S_i, r)$。引理 \ref{relative-cycles-lemma-descend-family} 中的
唯一性蕴含 $\alpha_i$ 的基变换是 $\alpha$，故 (1) 成立。
\end{proof}

\noindent
{\bf 讨论。}
若 $f : X \to S$、$r$ 与 $\alpha$ 如引理
\ref{relative-cycles-lemma-relative-r-cycle-general} 所述，则当引理
\ref{relative-cycles-lemma-relative-r-cycle-general} 的等价条件 (1) 与
(2) 成立时，可以称 $\alpha$ 为{\it $X/S$ 上的相对 $r$-循环}。这个定义有
许多良好性质；例如，当 $S$ 为 Noether 概形时，它与先前的定义不冲突，
且第 \ref{relative-cycles-section-families-specialization} 节的大多数结果都推广到这一
设定。

\medskip\noindent
还可作如下进一步推广。设 $S$ 任意，且 $f : X \to S$ 局部有限表示。设
$r \geq 0$，并令 $\alpha$ 为 $X/S$ 的纤维上的 $r$-循环族 $\alpha$。
如果对任意满足 $f(U) \subset V$ 的仿射开集 $U \subset X$ 与
$V \subset S$，限制 $\alpha|_U$ 都是上一段意义下 $U/V$ 上的相对
$r$-循环，则称 $\alpha$ 为{\it $X/S$ 上的相对 $r$-循环}。同样，先前的
许多结果都推广到这一设定。

\medskip\noindent
若以后需要这些推广，我们将在此仔细陈述并证明它们。
